% Standalone source generated from the accompanying modular TeX.
% Regenerate using python3 code/build_standalone.py.
% BEGIN preamble.tex
\documentclass[11pt,a4paper]{article}
\usepackage[T1]{fontenc}
\usepackage[utf8]{inputenc}
\usepackage{lmodern}
\usepackage[margin=25mm,headheight=14pt]{geometry}
\usepackage{microtype}
\usepackage{amsmath,amssymb,amsthm,mathtools,bm}
\usepackage{mathrsfs}
\usepackage{booktabs,longtable,array,enumitem}
\usepackage[dvipsnames]{xcolor}
\usepackage{fancyhdr}
\usepackage{hyperref}
\hypersetup{colorlinks=true,linkcolor=MidnightBlue,citecolor=MidnightBlue,
 urlcolor=RoyalBlue,pdftitle={Harmonic Shift Singularities and Exact Moment Identities},
 pdfauthor={ProveIt research continuation}}
\usepackage[nameinlink,noabbrev]{cleveref}
\numberwithin{equation}{section}
\newtheorem{theorem}{Theorem}[section]
\newtheorem{lemma}[theorem]{Lemma}
\newtheorem{proposition}[theorem]{Proposition}
\newtheorem{corollary}[theorem]{Corollary}
\theoremstyle{definition}
\newtheorem{definition}[theorem]{Definition}
\newtheorem{question}[theorem]{Question}
\theoremstyle{remark}
\newtheorem{remark}[theorem]{Remark}
\DeclareMathOperator{\Li}{Li}
\DeclareMathOperator{\FP}{FP}
\DeclareMathOperator{\CT}{CT}
\DeclareMathOperator{\Res}{Res}
\DeclareMathOperator{\rank}{rank}
\DeclareMathOperator{\Span}{span}
\newcommand{\dd}{\,\mathrm d}
\newcommand{\e}{\mathrm e}
\newcommand{\ii}{\mathrm i}
\newcommand{\src}[1]{\texttt{\detokenize{#1}}}
\newcommand{\C}{\mathbb C}
\newcommand{\Q}{\mathbb Q}
\newcommand{\N}{\mathbb N}
\newcommand{\Z}{\mathbb Z}
\newcommand{\Rea}{\operatorname{Re}}
\setlist{itemsep=3pt,topsep=5pt}
\setlength{\parindent}{1.3em}
\setlength{\parskip}{3pt}
\setlength{\emergencystretch}{2em}
\pagestyle{fancy}
\fancyhf{}
\fancyhead[L]{\small\sffamily ProveIt: exact harmonic and moment identities}
\fancyhead[R]{\small\sffamily Research continuation}
\fancyfoot[C]{\thepage}
\renewcommand{\headrulewidth}{0.3pt}

% END preamble.tex

\title{\textbf{Harmonic Shift Singularities and Exact Moment Identities}\\[5pt]
\large Spectral derivatives, polygamma moment spaces, harmonic powers,\\
and the finite constants of colliding periodic singularities}
\author{Research continuation for the ProveIt project\\
\small Prepared for Vladimir Reshetnikov}
\date{11 October 2026}
\begin{document}
\maketitle
% BEGIN sections/00_abstract.tex
\begin{abstract}
We continue the exact-identity programme of the ProveIt manuscript and
its incoming reports at a fixed repository revision. Four related
problems admit rigorous extensions. First, a Gamma-normalized harmonic
interpolation has an explicit singular germ at every nonpositive shift.
This determines its regular boundary constants in Arakawa--Kaneko
functions, its exact shift Taylor radii, and a functional rank jump:
at spectral order $-m$ the depth family has rank $m+1$, whereas any
positive spectral derivative restores infinite rank. Normalized local
primitives connect the same germs to the Lerch transcendent and
polylogarithms. Second, the mixed Bell rows generated by a Dougall
identity have exact rank $\lfloor w/2\rfloor$ in weight $w$; their
ordinary summation image has rank three for $w\ge6$. Explicit rational
elimination at weights six and eight yields a quartic polygamma
identity and an elementary generalized-harmonic tail-square evaluation.
Third, a triangular exponent generator proves that, for fixed $p$,
all centered values $F_p(-2r,1/2)$ lie in one finite rational span of
convergent Euler sums of weight at most $p-1$, independent of $r$.
We obtain complete power generators at $0,-2,-4$, all-moment Bernoulli
formulas for $p=1,2$, and exact exponent poles. A sufficient parity
class and finite ordinary-zeta formulas extend the calculation to
odd generalized harmonic orders. Fourth, a finite
principal-part formula for colliding periodic polygamma factors leaves
a jointly holomorphic remainder. It gives exact constants for nested
collision paths, arbitrary argument-derivative orders, and several
fully evaluated examples involving Stieltjes constants.

The proofs specify branches, continuation order, cutoff coordinates,
and the meaning of every rank assertion. Source specializations are
identified explicitly. Exact symbolic checks and independent numerical
diagnostics accompany the article; no numerical fit is used as a proof.
The current Gaussian $S_6$ and revised $S_8$ conjectures remain open.
We conclude with precise integration notes and further research targets.
\end{abstract}

\noindent\textbf{Keywords.}
Polylogarithm; Hurwitz zeta; generalized Stieltjes constant;
polygamma moment; harmonic power; analytic continuation;
Hadamard finite part; multiple zeta value.

\medskip
\noindent\textbf{Source baseline.}
ProveIt revision \src{921b58f10995}, frozen for this report on
11 October 2026. Full revision and archive provenance are recorded
in Section~\ref{sec:audit} and the accompanying integration notes.

% END sections/00_abstract.tex

\clearpage
\begingroup
\small
\setlength{\parskip}{0pt}
\tableofcontents
\endgroup
\clearpage
% BEGIN sections/01_scope.tex
\section{Scope, conventions, and the precise advances}
\label{sec:scope}

The aim is to turn several questions already present in the ProveIt
research stream into exact theorems. The most direct antecedents are
the mixed spectral report, the independent-order Hurwitz report, and
the raw harmonic-power section of the exact resonant report
\cite{RepoMixed,RepoIndependent,RepoExactResonant}. These sources
already contain substantial analytic continuation, Gamma normalization,
and finite-part machinery. We use it with explicit attribution, supply
self-contained proofs of the identities needed here, and resolve
specific questions about singularities, spans, and collision constants.

The literature background is classical: Gamma and polygamma
recurrences, Hurwitz-zeta continuation, Dougall's summation, and
polylogarithmic Mellin transforms
\cite{DLMFPolygamma,DLMFHurwitz,DLMFDougall,DLMFPolylog}.
The Arakawa--Kaneko functions are established objects
\cite{ArakawaKaneko,KanekoTsumura}. Our claims of an advance are
relative to the inspected ProveIt material. We do not claim that every
special-value consequence is absent from the wider literature.

\subsection{What is established}

\begin{longtable}{@{}p{0.20\textwidth}p{0.36\textwidth}p{0.38\textwidth}@{}}
\toprule
Problem & Antecedent & Result of this continuation\\
\midrule
\endhead
Shift singularities & A depth-dependent continuation half-plane;
the broader relative ordered Hurwitz character and its transport.
& Exact local amplitudes, boundary constants, sharp Taylor radii,
normalized primitives, and the finite-to-infinite functional rank
transition; Section~\ref{hsh:section}.\\[6pt]
Polygamma moments & A two-parameter Dougall identity and the request
to separate moments at weights six and eight.
& The full mixed-Bell rank in every weight, exact elimination
matrices, the three-dimensional summation image, and its full kernel;
Section~\ref{sec:bell}.\\[6pt]
Centered harmonic powers & Regularity of $F_p(-2r,1/2)$ and a
multiple-zeta expression obtained by finite word calculations.
& A triangular generator for every $r$ and all $p$, and one finite
Euler-sum space independent of $r$ at fixed $p$;
Section~\ref{hpnew:section}.\\[6pt]
Geometric collisions & Fixed-rate triple digamma formulas and
questions about nested rates and higher argument derivatives.
& An exact all-factor polygamma correction with jointly holomorphic
remainder, oriented harmonic hierarchy laws, and explicit derivative
examples; Section~\ref{hc:section}.\\
\bottomrule
\end{longtable}

The harmonic generator answers the all-power minus-two target and
gives every further even moment by finitely many triangular steps.
It is not a single closed two-variable expression in both the power
and the moment order. A companion result proves regularity for
polynomials in odd generalized harmonic orders and evaluates each
single odd-order family in a fixed finite space of even zeta values.
The collision theorem treats Stieltjes index
zero with arbitrary argument derivatives. Positive Stieltjes indices,
which introduce logarithmic local singularities, remain a separate
extension. The Bell rank concerns the stated formal polynomial rows;
it does not decide arithmetic independence of special values.
These distinctions delimit what has actually been proved.

\subsection{Functions and derivative conventions}

We use the ordinary Gamma function, $\psi=\Gamma'/\Gamma$, and
$\psi^{(q)}=\partial_a^q\psi$. The generalized harmonic numbers are
\[
 H_n^{(q)}(a)=\sum_{j=0}^{n-1}(a+j)^{-q},\qquad
 H_n^{(q)}=H_n^{(q)}(1),\qquad H_n=H_n^{(1)}.
\]
For an integer $q\ge1$ they satisfy
\begin{equation}\label{scope:harmonic-polygamma}
 H_n^{(q)}(a)=\frac{(-1)^{q-1}}{(q-1)!}
       \bigl\{\psi^{(q-1)}(a+n)-\psi^{(q-1)}(a)\bigr\}.
\end{equation}
The same formulas extend in $a$ away from their poles, with branches
specified when nonintegral powers occur. Relation
\eqref{scope:harmonic-polygamma} follows by telescoping the polygamma
recurrence \cite{DLMFPolygamma}.

For $\Re s>1$ and $a>0$,
$\zeta(s,a)=\sum_{n\ge0}(n+a)^{-s}$, and
$\zeta(s)=\zeta(s,1)$. Generalized Stieltjes constants have the
fixed convention
\begin{equation}\label{scope:Stieltjes}
 \zeta(1+v,a)=\frac1v+
       \sum_{j\ge0}\frac{(-1)^j}{j!}\gamma_j(a)v^j,
 \qquad \gamma_0(a)=-\psi(a),\qquad\gamma_j=\gamma_j(1).
\end{equation}
Thus $\gamma_j^{(q)}(a)$ always denotes an argument derivative,
whereas $\partial_s^k\zeta(s,a)$ is a spectral derivative.
The index $j$ in $\gamma_j$ is discrete and is never differentiated.
Similarly, a formula known at integer spectral labels $s=-m$ cannot
be differentiated with respect to $m$ to infer a spectral derivative.

Our multiple-polylogarithm convention is
\[
 \Li_{k_1,\ldots,k_d}(z)=
 \sum_{n_1>\cdots>n_d\ge1}
       \frac{z^{n_1}}{n_1^{k_1}\cdots n_d^{k_d}}.
\]
For admissible indices, its value at $z=1$ is
$\zeta(k_1,\ldots,k_d)$; weight means $k_1+\cdots+k_d$.
The one-index case is the ordinary polylogarithm
\cite{DLMFPolylog}. The Lerch transcendent is initially
\[
 \Phi(z,s,a)=\sum_{n\ge0}\frac{z^n}{(n+a)^s},\qquad |z|<1,
\]
with analytic continuation where required
\cite{DLMFLerch}. In particular $z\Phi(z,s,1)=\Li_s(z)$.

\subsection{Three different meanings of a displayed infinite expression}

Most sums in Section~\ref{sec:bell}, including the tail-square
identity, converge absolutely in the ordinary sense. They need no
regularization. In contrast,
\[
 F_p(-2r,1/2)
\]
means the value of the meromorphic continuation of
$\sum_{n\ge0}H_n^p(n+1/2)^{-s}$ at the regular point $s=-2r$.
For $p\ge1$ this is not the sum of an ordinarily convergent numerical
series after direct substitution of $s=-2r$.

For a function with a local Laurent--log expansion, $\CT_z$ extracts
the coefficient of $z^0(\log z)^0$. A Hadamard finite-part integral
is the constant term in its cutoff expansion with the cutoff coordinate
held fixed. At distinct interior seams we use independent ambient
coordinates $x-x_i$ and take their joint constant term before a
collision parameter tends to zero. The subsequent $\CT_\epsilon$
in Section~\ref{hc:section} is a second operation, on the resulting
function of the geometric shifts. Changing either coordinate can
change this constant. The exact correction formulas record that
dependence.

Branches are also part of the statements. A principal shift
continuation, a continuation around a negative integer, and a real
collision path need not have the same analytic germ. All local powers
are written using an explicitly chosen logarithm. Fixed-endpoint
primitives are normalized to vanish at their stated base point, so
an integration constant cannot absorb a resonance discrepancy.

\subsection{Proof strategy and verification}

Each family is handled by first retaining a full analytic parameter
family, then making the desired specialization. Finite transport
isolates a shift singularity; a Gamma quotient generates the Bell
moments; an exponent parameter linearizes raw harmonic powers; and
partial fractions isolate the singular collision pieces. Normal
convergence or a holomorphic remainder is proved before extracting
coefficients. This order of operations is essential at resonances.

Exact symbolic checks replay coefficient identities and rational
linear algebra. Independent numerical evaluations start from different
representations and diagnose signs, constants, branches, and omitted
terms. They supplement the proofs. The numerical tails and quadratures
are not interval certificates, and their small residuals are not used
to justify an infinite identity. Section~\ref{sec:audit} records the
actual checks and the boundaries of the source audit.

% END sections/01_scope.tex

% BEGIN sections/02_shift_germs.tex
\section{Exact shift singularities of the entire harmonic difference}
\label{hsh:section}

The Newton expansion in the mixed-spectral report \cite{RepoMixed}
proves holomorphy of its depth-$r$ entire difference on $\Re a>-r$,
and asks for the actual singularity at $a=-r$.  We determine every
negative-integral shift germ, the exact Taylor radius, and the finite
regular term at the first singularity.  The proof uses a finite shift
identity and a cancellation at $a=1$.  Finite transport is a classical
special-function method.  The independent-orders report
\cite{RepoIndependent} supplies a broader entire ordered interpolation
of which this remainder is an exact specialization, identified below.
The contribution here is the explicit singular amplitudes and their
boundary, Taylor-radius, and functional-rank consequences; we do not
assert priority for the underlying transport method.

\subsection{Definitions and the finite shift identity}

For $\Re a>0$, define, initially for $\Re s>1+\Re u$,
\begin{align}
 F_a(s,u)&=\sum_{n\geq0}\frac{(a+u)_n}{(a)_n(n+a)^s},
 &K_a(u)&=\frac{\Gamma(a)\Gamma(1+u)}{\Gamma(a+u)},\label{hsh:definitions}\\
 D(s,u;a)&=F_a(s,u)-K_a(u)F_1(s,u),
 &D_r(s;a)&=[u^r]D(s,u;a).\label{hsh:difference}
\end{align}
Here $(x)_n=x(x+1)\cdots(x+n-1)$ is a rising factorial.  Every
statement in $u$ can be read on a fixed disk $|u|<\rho<1$.
The sources \cite{RepoExact,RepoMixed} prove that $D$ is entire in
$s$, and that its coefficients are jointly holomorphic for
$s\in\mathbb C$, $\Re a>-r$.  In particular,
\begin{equation}\label{hsh:basezero}
 D(s,u;1)=0,\qquad D_0(s;a)=\zeta(s,a)-\zeta(s).
\end{equation}
The pole of the last difference at $s=1$ is removable.

Put
\begin{equation}\label{hsh:P}
 P_n(a,u)=\frac{(a+u)_n}{(a)_n}
       =\sum_{j=0}^n e_j(n;a)u^j,
 \qquad
 e_j(n;a)=\sum_{0\leq\nu_1<\cdots<\nu_j<n}
                  \prod_{i=1}^j\frac1{a+\nu_i}.
\end{equation}
The empty product is one and $e_j(n;a)=0$ for $j>n$.

\paragraph{Exact identification with the ordered endpoint interpolation.}
Let $\mathcal H_{b,a}(w)$ denote the relative character in
\cite{RepoIndependent}, normalized by the deconcatenation identity
$Z_a(w)=\sum_{w=vw'}Z_b(v)\mathcal H_{b,a}(w')$ for the
strict ordered Hurwitz sums.  Its Gamma formula for words of ones is
$\mathcal H_{1,a}(\{1\}^j)=[u^j]K_a(u)$.  Applying the defining
recursion to $(s,\{1\}^r)$ therefore gives the exact identification
\begin{equation}\label{hsh:relativecharacter}
          D_r(s;a)=\mathcal H_{1,a}(s,\{1\}^r).
\end{equation}
Indeed $Z_a(s,\{1\}^r)=\sum_{n\geq0}e_r(n;a)(n+a)^{-s}$,
and the other deconcatenation terms are exactly the coefficients
of $K_a(u)F_1(s,u)$ in \eqref{hsh:difference}.  This first proves
the identity in a convergence domain, then everywhere by the entire
spectral continuations.  Proposition~\ref{hsh:transport} is accordingly
also a specialization of that source's endpoint transport.

\begin{proposition}[Exact finite transport]\label{hsh:transport}
For every integer $L\geq1$,
\begin{equation}\label{hsh:finite}
 D(s,u;a)=\sum_{n=0}^{L-1}P_n(a,u)(a+n)^{-s}
                       +P_L(a,u)D(s,u;a+L).
\end{equation}
Consequently
\begin{align}
 D(s,u;a)&=a^{-s}+\frac{a+u}{a}D(s,u;a+1),
                                                     \label{hsh:one}\\
 D_r(s;a)&=\mathbf1_{r=0}a^{-s}+D_r(s;a+1)
                      +\mathbf1_{r\geq1}\frac{D_{r-1}(s;a+1)}a,
                                                     \label{hsh:depthshift}\\
 D_r(s;a)&=\sum_{n=r}^{L-1}e_r(n;a)(a+n)^{-s}
          +\sum_{j=0}^{\min(r,L)}e_j(L;a)D_{r-j}(s;a+L).
                                                     \label{hsh:depthfinite}
\end{align}
The finite sum in the last formula is empty when $L\leq r$.
These identities hold wherever the powers and continuations on both
sides use the same branches.
\end{proposition}

\begin{proof}
Separating the $n=0$ summand gives
$F_a=a^{-s}+(a+u)F_{a+1}/a$.  The Gamma functional equation gives
$K_a=(a+u)K_{a+1}/a$.  Subtraction proves
\eqref{hsh:one} in the initial common half-plane.  Both sides are
entire in $s$, so that identity continues to every $s$.  Iterating it
gives \eqref{hsh:finite}; coefficient extraction gives the other
formulas.  All operations are finite after the initial identity, so
they create no interchange of a divergent sum and a derivative.
\end{proof}

\subsection{Every negative-integral shift germ}

For $N\geq0$, put $a=-N+z$ and define
\begin{equation}\label{hsh:C}
 C_N(z,u)=P_N(-N+z,u)
          =\prod_{j=1}^{N}\left(1+\frac{u}{z-j}\right),
 \qquad c_{N,r}(z)=[u^r]C_N(z,u).
\end{equation}
Thus $C_0=1$ and $c_{N,r}=0$ if $r>N$.  Near $z=0$ choose
a branch of $\log z$.  For the other powers in
\eqref{hsh:finite}, choose logarithms of $z-h$, $1\leq h\leq N$,
holomorphic on a small disk around zero.  Their values at zero may be
$\log h+\mathrm i\pi$ or $\log h-\mathrm i\pi$ for the two lateral
continuations.

\begin{theorem}[Exact local decomposition]\label{hsh:germ}
On these specified branches,
\begin{equation}\label{hsh:local}
 D(s,u;-N+z)=C_N(z,u)z^{-s}+H_N(s,u;z),
\end{equation}
where $H_N$ is jointly holomorphic in $s$, $u$, and $z$ near
every point $(s_0,0,0)$.  An explicit formula is
\begin{align}
 H_N(s,u;z)={}&
 \sum_{n=0}^{N-1}P_n(-N+z,u)(z-N+n)^{-s}\notag\\
 &+\left(1+\frac uz\right)C_N(z,u)D(s,u;1+z).
                                                     \label{hsh:H}
\end{align}
The apparent pole in the second line is removable.  For every
$r,k\geq0$,
\begin{equation}\label{hsh:localjets}
 \partial_s^kD_r(s;-N+z)
   =c_{N,r}(z)z^{-s}(-\log z)^k
                         +\partial_s^kH_{N,r}(s;z),
 \qquad H_{N,r}=[u^r]H_N.
\end{equation}
For $r\leq N$ the leading coefficient is nonzero:
\begin{equation}\label{hsh:leading}
 c_{N,r}(0)=(-1)^r e_r\left(1,\frac12,\ldots,\frac1N\right).
\end{equation}
For $r>N$ the entire singular amplitude vanishes identically.
\end{theorem}

\begin{proof}
Use \eqref{hsh:finite} with $L=N+1$.  Its $n=N$ term is exactly
$C_N(z,u)z^{-s}$, and the other terms give \eqref{hsh:H} because
$P_{N+1}(-N+z,u)=(1+u/z)C_N(z,u)$.
Every factor of $P_n$, $n\leq N$, is holomorphic in $z$ near
zero.  The chosen logarithms make all powers with $n<N$
holomorphic there as well.  By \eqref{hsh:basezero},
$D(s,u;1+z)=z\,\widetilde D(s,u;z)$ with $\widetilde D$
jointly holomorphic.  This cancels the sole $1/z$ factor in the
last line.  This also proves local normality on compact spectral
sets, so all the stated derivatives can be taken.  Finally,
$C_N(0,u)=\prod_{j=1}^{N}(1-u/j)$, which proves
\eqref{hsh:leading}; its elementary symmetric coefficients are
strictly positive before the displayed sign.
\end{proof}

\begin{corollary}[Principal continuation]\label{hsh:principal}
For each $r\geq0$, $D_r(s;a)$ has a jointly holomorphic principal
continuation to
\begin{equation}\label{hsh:slit}
       (s,a)\in\mathbb C\times
                         \bigl(\mathbb C\setminus(-\infty,-r]\bigr).
\end{equation}
It is holomorphic across $a=0,-1,\ldots,-(r-1)$.  On the
specified lateral branches, every $a=-N$ with $N\geq r$ has
the germ \eqref{hsh:localjets}.
\end{corollary}

\begin{proof}
On any compact set in the slit plane, choose $L$ so large that
$\Re(a+L)>0$.  In \eqref{hsh:depthfinite}, the only powers that
remain have $n\geq r$, hence their arguments avoid the principal
logarithm cut.  The tail is defined in the original half-plane.
The rational coefficients have only possible poles at negative
integers.  Those in the slit domain have $N<r$, and
Theorem~\ref{hsh:germ} proves their removability.  The formulas
agree on overlaps by finite transport and coincide with the
original function on a common open set.  They therefore give
the claimed jointly holomorphic continuation.
\end{proof}

\paragraph{A distinction between principal continuation and all sheets.}
The early removable points cannot simply be deleted from the
puncture set for every sheet.  A positive circuit around $a=-N$
changes the chosen germ by
\begin{equation}\label{hsh:monodromy}
 \Delta_N D_r(s;a)=
 (\exp(-2\pi\mathrm i s)-1)e_r(N;a)(a+N)^{-s}.
\end{equation}
This follows by changing only $\log(a+N)$ in a finite transport
formula with $L>N$.  The rational coefficient $e_r(N;a)$ may have
poles at earlier shifts.  For instance, a circuit around $-1$
changes $D_1(s;a)$ by
\[
       (\exp(-2\pi\mathrm i s)-1)\frac{(a+1)^{-s}}a.
\]
Although the principal $D_1$ is regular at $a=0$, this continued
branch has a pole there for generic $s$.  At $s=-m$, differentiating
once in $s$ gives the change $-2\pi\mathrm i(a+1)^m/a$.
Thus the same issue occurs for spectral derivatives at the polynomial
orders.  Formula \eqref{hsh:slit} is a statement about the principal
branch.  Arbitrary multivalued continuation can safely be made along
paths avoiding all of $0,-1,-2,\ldots$; it need not preserve the
removability of every preimage of an earlier shift.  This is a
continuation qualification, not a counterexample to the source's
half-plane theorem.

\subsection{The first boundary and Arakawa--Kaneko values}

For $j\geq1$, write
\begin{equation}\label{hsh:Xi}
 \Xi_j(s)=\frac1{\Gamma(s)}\int_0^\infty
   \frac{t^{s-1}}{e^t-1}\operatorname{Li}_j(1-e^{-t})\,dt,
 \qquad \Re s>0.
\end{equation}
These are the classical Arakawa--Kaneko functions, continued entirely
in $s$ by Taylor subtraction at zero
\cite{ArakawaKaneko,KanekoTsumura}.  In the descending convention,
the multiple version $\Xi(j,\{1\}^{p-1};s)$ is obtained by replacing
the polylogarithm in \eqref{hsh:Xi} by
$\operatorname{Li}_{j,1,\ldots,1}$.
The mixed-spectral report proves the exact shift-jet identity
\begin{equation}\label{hsh:AKjet}
 \left.\partial_a^pD_q(s;a)\right|_{a=1}
      =(-1)^{q+p}p!\,\Xi(q+1,\{1\}^{p-1};s),\qquad p\geq1.
\end{equation}
In particular $\Xi_1(s)=s\zeta(s+1)$, with its removable value
at $s=0$ understood.  This follows directly from
$\operatorname{Li}_1(1-e^{-t})=t$.

For a fixed positive depth $r$, set
\begin{align}
 E_h^{(r)}&=e_h\left(1,\frac12,\ldots,\frac1r\right),
 \qquad\prod_{j=1}^r\left(1-\frac zj\right)^{-1}
       =\sum_{n\geq0}h_n^{(r)}z^n,\label{hsh:eh}\\
 c_r(z)&=c_{r,r}(z)=\frac1{\prod_{j=1}^r(z-j)}
       =\frac{(-1)^r}{r!}\prod_{j=1}^r\left(1-\frac zj\right)^{-1},\notag\\
 B_r(s)&=(-1)^r\sum_{h=0}^{r-1}E_h^{(r)}\Xi_{r-h}(s).
                                                       \label{hsh:B}
\end{align}
The $h_n^{(r)}$ are complete homogeneous symmetric polynomials.

\begin{theorem}[Exact first-boundary normalization]\label{hsh:boundary}
For $r\geq1$,
\begin{equation}\label{hsh:boundaryidentity}
 \boxed{\displaystyle
 \lim_{z\to0}\left\{D_r(s;-r+z)-c_r(z)z^{-s}\right\}=B_r(s).}
\end{equation}
The expression inside braces extends holomorphically through $z=0$.
Convergence is locally uniform in $s$, and arbitrary fixed spectral
derivatives can be applied to the identity.  In particular,
\begin{align}
 B_1(s)&=-s\zeta(s+1),\notag\\
 B_2(s)&=\Xi_2(s)+\tfrac32s\zeta(s+1),\notag\\
 B_3(s)&=-\Xi_3(s)-\tfrac{11}{6}\Xi_2(s)-s\zeta(s+1).
                                                       \label{hsh:examples}
\end{align}
Every Taylor coefficient of the holomorphic remainder in
\eqref{hsh:boundaryidentity} is a finite combination of the multiple
functions in \eqref{hsh:AKjet}, with rational coefficients.
\end{theorem}

\begin{proof}
Take $N=r$ in \eqref{hsh:H} and extract $u^r$.  The finite terms
with $n<r$ vanish because $P_n$ has degree $n$.  Thus the remainder
is exactly
\begin{equation}\label{hsh:firstH}
 H_{r,r}(s;z)=[u^r]\left(1+\frac uz\right)
                          C_r(z,u)D(s,u;1+z).
\end{equation}
At $z=0$, the term without $u/z$ vanishes.  The remaining term is
$[u^r]uC_r(0,u)\partial_aD(s,u;1)$.
Now $C_r(0,u)=\sum_{h=0}^r(-1)^hE_h^{(r)}u^h$, while
\eqref{hsh:AKjet} at $p=1$ gives
$\partial_aD(s,u;1)=\sum_{q\geq0}(-1)^{q+1}\Xi_{q+1}(s)u^q$.
Their finite coefficient product is precisely $B_r$.
The holomorphy assertion is already contained in
Theorem~\ref{hsh:germ}.  For higher Taylor coefficients, expand
$C_r(z,u)$ rationally at $z=0$ and use
\[
 D(s,u;1+z)=\sum_{p\geq1}\sum_{q\geq0}
       (-1)^{p+q}\Xi(q+1,\{1\}^{p-1};s)u^qz^p.
\]
At a fixed power of $z$ in \eqref{hsh:firstH}, only finitely many
$p,q$ occur.  This also proves the final assertion with its
normalization specified.
\end{proof}

The same calculation gives every regular boundary constant.  With
$E_h^{(N)}=e_h(1,1/2,\ldots,1/N)$ and $E_0^{(0)}=1$,
\begin{equation}\label{hsh:allregular}
 \begin{split}
 H_{N,r}(s;0)={}&\sum_{n=0}^{N-1}e_r(n;-N)(n-N)^{-s}\\
 &+(-1)^r\sum_{h=0}^{\min(N,r-1)}E_h^{(N)}\Xi_{r-h}(s).
 \end{split}
\end{equation}
The second sum is empty for $r=0$.  This follows by evaluating the
first line of \eqref{hsh:H} at zero and extracting the coefficient
of $u^r$ from $uC_N(0,u)\partial_aD(s,u;1)$ in its second line.
The finite powers retain their specified lateral logarithms.  If
$r>N$, the first sum is empty and the singular amplitude vanishes;
the formula recovers the source's previously proved negative-shift
values.  If $r=N\geq1$, it is \eqref{hsh:boundaryidentity}.

\paragraph{The regular boundary term and a coordinate finite part.}
The subtraction in \eqref{hsh:boundaryidentity} retains the full
analytic amplitude $c_r(z)$; it is not the same operation as discarding
only the divergent terms of a Laurent expansion.  For $p\geq0$,
take the finite part along the positive real coordinate $z$ to mean
the coefficient of $z^0(\log z)^0$ in the power--logarithm expansion.
Then
\begin{equation}\label{hsh:FP}
 \boxed{\displaystyle
 \FP_{z=0}D_r(p;-r+z)=B_r(p)+\frac{(-1)^r}{r!}h_p^{(r)}.}
\end{equation}
Indeed the constant contributed by $c_r(z)z^{-p}$ is
$[z^p]c_r(z)$.  On the other hand, for every $k\geq1$,
\begin{equation}\label{hsh:FPjets}
 \FP_{z=0}\left.\partial_s^kD_r(s;-r+z)\right|_{s=p}
                       =B_r^{(k)}(p),\qquad p\geq0,
\end{equation}
because its singular amplitude has the factor $(-\log z)^k$.
These formulas concern a fixed coordinate and a defined constant
coefficient, so there is no implicit conversion between two
regularization prescriptions.

For $m\geq1$ and $k\geq0$, there is an ordinary lateral limit,
\begin{equation}\label{hsh:negativeboundary}
 \lim_{z\to0}\left.\partial_s^kD_r(s;-r+z)\right|_{s=-m}
                      =B_r^{(k)}(-m),
\end{equation}
provided $\arg z$ stays in a bounded interval on the chosen branch.
The reason is $z^m(\log z)^k\to0$; no finite part is needed.
For positive $k$, this continuous extension is not
holomorphic, as the next subsection proves.

\paragraph{Explicit Stieltjes and Gamma boundary identities at depth one.}
Use the convention
$\zeta(1+s)=s^{-1}+\sum_{j\geq0}(-1)^j\gamma_js^j/j!$.
Then \cite{DLMFHurwitz}
\begin{equation}\label{hsh:Stieltjes}
 B_1(0)=-1,\qquad
 B_1^{(k)}(0)=(-1)^k k\gamma_{k-1}\quad(k\geq1).
\end{equation}
Consequently, for example,
\begin{align}
 \lim_{z\to0}\{D_1(0;-1+z)-(z-1)^{-1}\}&=-1,\notag\\
 \lim_{z\to0}\{\partial_sD_1(0;-1+z)
                          +(z-1)^{-1}\log z\}&=-\gamma,\notag\\
 \lim_{z\to0}\{\partial_s^2D_1(0;-1+z)
                          -(z-1)^{-1}(\log z)^2\}&=2\gamma_1.
                                                       \label{hsh:stieltjeslimits}
\end{align}
At the negative orders, for $m\geq1$,
\begin{equation}\label{hsh:zetaexplicit}
 B_1^{(k)}(-m)=m\zeta^{(k)}(1-m)
                         -k\zeta^{(k-1)}(1-m),
\end{equation}
where the last term is omitted for $k=0$.  In particular,
\[
 B_1'(-1)=\frac12-\frac12\log(2\pi),\qquad
 B_1'(-2)=\frac1{12}+2\zeta'(-1).
\]
The first evaluation uses the classical Lerch normalization of
$\zeta'(0)$ \cite{DLMFHurwitz}.  These constants occur here as
boundary values of a continued harmonic family, with the complete
singular amplitude specified.

\subsection{Normalized Lerch and polylogarithmic primitives}
\label{hsh:primitives}

The exact singular amplitude also gives a finite antiderivative formula.
This is an application of standard partial fractions and Lerch calculus
\cite{DLMFLerch}, with the normalization needed at resonant spectral
orders made explicit.  Fix a logarithm on a slit punctured unit disk
and a base point $z_0$ in that disk.  For $1\leq j\leq r$, put
\begin{align}
 A_{r,j}&=\frac{(-1)^{r-j}}{(j-1)!(r-j)!},
 &c_r(z)&=\sum_{j=1}^{r}\frac{A_{r,j}}{z-j},
                                                    \label{hsh:primitivePF}\\
 P_j(s;z,z_0)&=-\frac1j\left\{
 z^{1-s}\Phi(z/j,1,1-s)
       -z_0^{1-s}\Phi(z_0/j,1,1-s)\right\}.
                                                    \label{hsh:Lerchprimitive}
\end{align}
Here $\Phi(w,1,v)=\sum_{n\geq0}w^n/(n+v)$ for $|w|<1$,
continued meromorphically in $v$.

\begin{proposition}[Entire normalized spectral primitive]
\label{hsh:primitive}
The difference $P_j$ extends to an entire function of $s$, including
the apparent resonances $s=1,2,\ldots$, and
\begin{equation}\label{hsh:primitivederivative}
 P_j(s;z_0,z_0)=0,\qquad
 \frac{\partial}{\partial z}P_j(s;z,z_0)=\frac{z^{-s}}{z-j}.
\end{equation}
Consequently a normalized primitive of the complete first-boundary germ is
\begin{equation}\label{hsh:fullprimitive}
 \sum_{j=1}^{r}A_{r,j}P_j(s;z,z_0)
                   +\int_{z_0}^{z}H_{r,r}(s;t)\,\mathrm dt.
\end{equation}
The remainder in this formula is holomorphic on the unit disk.
Every fixed spectral derivative can be taken in the formula.
At $s=0$, all the differentiated singular primitives are finite
ordinary-polylogarithm expressions: define
\begin{equation}\label{hsh:polylogprimitive}
 J_{j,k}(z)=-k!\sum_{h=0}^{k}\frac{(-\log z)^h}{h!}
                          \operatorname{Li}_{k+1-h}(z/j).
\end{equation}
Then
\begin{equation}\label{hsh:polylogprimitivejet}
 \left.\partial_s^kP_j(s;z,z_0)\right|_{s=0}
                         =J_{j,k}(z)-J_{j,k}(z_0),
 \qquad J_{j,k}'(z)=\frac{(-\log z)^k}{z-j}.
\end{equation}
\end{proposition}

\begin{proof}
The residues of $c_r$ at its distinct simple poles give
\eqref{hsh:primitivePF}.  Expanding the Lerch functions gives the
normally convergent, already normalized series
\begin{equation}\label{hsh:primitiveentire}
 P_j(s;z,z_0)=
 -\sum_{m\geq1}\frac{z^{m-s}-z_0^{m-s}}{j^m(m-s)}.
\end{equation}
For a term with $s=m$, interpret the quotient as
$\log z-\log z_0$.  Each term is entire in $s$.  On compact
spectral sets and compact subsets of the slit punctured disk, its tail
is bounded by a geometric series times $1/m$.  Hence the sum and all
fixed spectral derivatives are locally normal.  Equivalently, each
unnormalized Lerch primitive in \eqref{hsh:Lerchprimitive} has residue
$j^{-p}$ at $s=p\geq1$, independent of its endpoint; the difference
cancels that residue.  Differentiating \eqref{hsh:primitiveentire}
in $z$ proves \eqref{hsh:primitivederivative}.  Formula
\eqref{hsh:firstH} shows that $H_{r,r}$ is holomorphic for $|z|<1$,
with its sole apparent pole at zero removed.  Thus its integral is
well defined independently of the path in that disk, and
\eqref{hsh:fullprimitive} follows.  Finally, differentiating
$z^{-s}/(m-s)$ at $s=0$ and summing the resulting powers $m^{-q}$
gives \eqref{hsh:polylogprimitive}; termwise differentiation, or the
usual polylogarithm derivative rule, proves
\eqref{hsh:polylogprimitivejet}.
\end{proof}

In particular,
\[
 J_{j,0}(z)=\log(1-z/j),\qquad
 J_{j,1}(z)=-\operatorname{Li}_2(z/j)
                        -\log z\log(1-z/j).
\]
The sign in the second formula is fixed by
$J_{j,1}'(z)=-\log z/(z-j)$.  The base-point subtraction supplies
the additive constant before continuation or spectral differentiation;
separate unnormalized Lerch terms should not be evaluated at their
resonant parameter poles.

\subsection{Exact Taylor radii and the rank jump at integral orders}

\begin{theorem}[Exact shift Taylor radius]\label{hsh:radius}
Fix a real center $a_*>-r$ and a spectral order $s_0$.
The Taylor series in $a-a_*$ of $D_r(s_0;a)$ has radius
\[
 \begin{cases}
 a_*+r,&s_0\notin\{0,-1,-2,\ldots\},\\
 \infty,&s_0\in\{0,-1,-2,\ldots\}.
 \end{cases}
\]
For every $k\geq1$ and every $s_0\in\mathbb C$, the Taylor
series of $\partial_s^kD_r(s_0;a)$ has radius exactly $a_*+r$.
In particular, at the center $a_*=1$, the lower bound $r+1$
in \cite{RepoMixed} is exact outside the polynomial cases, and is
exact for every positive spectral derivative even at those cases.
\end{theorem}

\begin{proof}
The disk $|a-a_*|<a_*+r$ lies in the domain
\eqref{hsh:slit}, so holomorphy gives the lower bound, uniformly
on compact spectral sets.  At $a=-r$, the singular amplitude is
$c_r(z)z^{-s_0}(-\log z)^k$, with $c_r(0)=(-1)^r/r!\ne0$;
for $r=0$ the same statement has $c_0=1$.
When $k=0$ and $s_0$ is not a nonpositive integer, $z^{-s_0}$
has either a pole or nontrivial branching.  It therefore cannot
extend holomorphically to zero.  For $k\geq1$, its logarithmic
factor prevents holomorphic extension for every $s_0$.  More
explicitly, one circuit sends the germ to
$e^{-2\pi\mathrm i s_0}z^{-s_0}(-\log z-2\pi\mathrm i)^k$.
If $s_0$ is nonintegral, the exponential factor is nontrivial;
if $s_0$ is integral, the nonconstant logarithmic polynomial is
not translation invariant.  Positive integral $s_0$ at $k=0$
instead gives a pole.  The holomorphic remainder cannot cancel
any of these obstructions.  Hence the Taylor disk cannot be
larger.  At $s_0=-m$, the finite Stirling formula below proves
that $D_r(-m;a)$ is a polynomial and has infinite radius.
\end{proof}

\begin{corollary}[Functional independence of spectral jets]
\label{hsh:independence}
If $s_0\notin\{0,-1,-2,\ldots\}$, every finite family
\[
 \{\partial_s^kD_r(s_0;a):0\leq r\leq R,\ 0\leq k\leq K\}
\]
is linearly independent over $\mathbb C$, even modulo entire
functions of $a$.  If $s_0=-m$ with $m\geq0$, the same assertion
holds for the subfamily $1\leq k\leq K$.
This is independence of functions of the shift, not arithmetic
independence of any special values.
\end{corollary}

\begin{proof}
Suppose a finite linear combination equals an entire function,
and let $r$ be the smallest depth with a nonzero coefficient.
Every larger-depth term is holomorphic at $a=-r$ on the principal
continuation.  The singular germ of the depth-$r$ terms is
$c_r(z)z^{-s_0}P(-\log z)$ for a nonzero polynomial $P$.
The monodromy argument from Theorem~\ref{hsh:radius} shows that
this cannot be holomorphic, except possibly when $s_0$ is a
nonpositive integer and $P$ is constant.  That exception is
excluded by the stated subfamily.  Thus all coefficients at
this depth vanish, a contradiction.  Induction gives the claim.
\end{proof}

Let $\left\{\begin{smallmatrix}n\\j\end{smallmatrix}\right\}$
denote a Stirling number of the second kind and let
$(a-1)^{\underline j}=(a-1)\cdots(a-j)$.
The known finite formula \cite{RepoExact,RepoMixed} is
\begin{equation}\label{hsh:finiteStirling}
 D_r(-m;a)=\sum_{j=1}^{m+1}
 \left\{\begin{matrix}m+1\\j\end{matrix}\right\}
           (a-1)^{\underline j}\left(-\frac1j\right)^{r+1}.
\end{equation}

\begin{proposition}[Exact finite rank and minimal depth recurrence]
\label{hsh:rank}
For $m\geq0$, the functions $\{D_r(-m;a):r\geq0\}$ span
exactly the $(m+1)$-dimensional space of polynomials of degree
at most $m+1$ vanishing at $a=1$.  The first $m+1$ depths form
a basis.  If $E$ shifts depth, $ED_r=D_{r+1}$, then
\begin{equation}\label{hsh:recurrence}
 \boxed{\displaystyle
       \prod_{j=1}^{m+1}\left(E+\frac1j\right)D_r(-m;a)=0,
       \qquad r\geq0.}
\end{equation}
Its order $m+1$ is minimal among nonzero constant-coefficient
depth recurrences valid as identities in $a$.
For any fixed $k\geq1$, the family
$\{\partial_s^kD_r(-m;a):r\geq0\}$ is infinite dimensional and
satisfies no nonzero constant-coefficient depth recurrence.
The latter statement also holds at nonexceptional $s_0$ with $k=0$.
\end{proposition}

\begin{proof}
The polynomials
$\left\{\begin{smallmatrix}m+1\\j\end{smallmatrix}\right\}
(a-1)^{\underline j}$, $1\leq j\leq m+1$, have distinct degrees,
nonzero leading coefficients, and vanish at one.  They form a basis
of the stated space.  The coefficient matrix for depths $0,\ldots,m$
in \eqref{hsh:finiteStirling} has entries $\lambda_j^{r+1}$,
where $\lambda_j=-1/j$.  Its determinant is a nonzero product
of the $\lambda_j$ and a Vandermonde determinant.  This proves
the rank and basis assertions.  Each component in
\eqref{hsh:finiteStirling} is annihilated by $E-\lambda_j$,
which proves \eqref{hsh:recurrence}.  Conversely, independence
of those polynomial components forces an annihilating polynomial
$Q$ to vanish at every $\lambda_j$.  Its degree is therefore
at least $m+1$.  Finally any recurrence for a positive spectral
derivative would give a finite constant linear relation among
distinct depths, contradicting Corollary~\ref{hsh:independence}.
\end{proof}

For example, at $s=0$ one has
$D_r(0;a)=(-1)^{r+1}(a-1)$ and $(E+1)D_r=0$.
At $s=-1$ the minimal recurrence is
$D_{r+2}+\tfrac32D_{r+1}+\tfrac12D_r=0$.
A single spectral differentiation removes these finite-rank
relations.  This failure is detected by exact logarithmic germs,
without any numerical independence assumption.

\subsection{Audit and reproducibility}

The source's holomorphy domain and Taylor lower bound were correct;
the new statements determine their actual boundary behavior.  The
regular boundary term $B_r$ and the coordinate finite part in
\eqref{hsh:FP} must retain their different definitions.  Likewise,
principal-sheet removability must not be promoted to removability on
all sheets: \eqref{hsh:monodromy} supplies an explicit obstruction.

The accompanying program \src{code/verify_shift_germs.py} checks the
finite recurrence and residue algebra exactly with rational arithmetic,
compares the Newton series with the independently convergent defining
harmonic sums at complex parameters, checks finite transport at negative
shifts, and compares Cauchy-extracted boundary constants with
Arakawa--Kaneko integrals and ordinary zeta derivatives.  The numerical
checks are finite-precision diagnostics with separately recorded
truncations; they are not interval certificates or replacements for
the normal-convergence and continuation arguments above.
The short companion \src{code/verify_shift_primitives.py} checks the
partial-fraction amplitudes and the polylogarithmic primitive derivatives
with exact symbolic arithmetic.

% END sections/02_shift_germs.tex

% BEGIN sections/03_bell_moments.tex
\section{The exact space of mixed polygamma moments}
\label{sec:bell}

The mixed-spectral report asks which individual polygamma moments can
be separated from its two-parameter Gamma identity, with explicit first
targets at total orders six and eight \cite{RepoMixed}. This section
answers that question for the stated family at every weight. There are
two different dimensions to compute: the dimension of its formal
polynomial rows, and the dimension of the functions obtained after
summation. The former grows with the weight; the latter is exactly three
from weight six onward. Their difference produces an infinite family of
identically vanishing, absolutely convergent polygamma sums.

\subsection{Normalization and the underlying Gamma identity}

For $\Re a>0$, set $x_n=n+a/2$, $A=a-1$, and
\begin{equation}
 L_k(n;a)=(-1)^k\psi^{(k-1)}(n+a)-\psi^{(k-1)}(n+1),
 \qquad k\ge2.
 \label{bell:lambda}
\end{equation}
These are ordinary functions, and all products below are algebraic
products, without complex conjugation. Write
\begin{align}
 Q_n(p,q;a)
 &=\frac{\Gamma(n+a)\Gamma(n+1+p)\Gamma(n+1+q)
              \Gamma(n+a-p-q)}
 {\Gamma(n+1)\Gamma(n+a-p)\Gamma(n+a-q)
              \Gamma(n+1+p+q)},\label{bell:Q}\\
 b_a(p,q)
 &=\psi(a/2)+\tfrac12\{\psi(2)-\psi(1+p)-\psi(1+q)
                                      -\psi(a-p-q)\}.
 \label{bell:b}
\end{align}

\begin{lemma}[The double-zero Gamma sum: antecedent]
\label{bell:antecedent}
For sufficiently small $p,q$, locally uniformly with $\Re a>0$,
\begin{equation}
 \sum_{n\ge0}\left\{x_n(Q_n(p,q;a)-1)
             +\frac{(A-p-q)pq}{x_n}\right\}
             =-(A-p-q)pq\,b_a(p,q).
 \label{bell:master}
\end{equation}
Every fixed parameter derivative can be taken in the displayed bracket.
\end{lemma}

\begin{proof}
This is the double-zero identity proved in \cite{RepoMixed}. We give
the residue calculation to fix its normalization. In the Rogers--Dougall
${}_5F_4(1)$ sum \cite[16.4.9]{DLMFDougall}, set
$b=1+p$, $c=1+q$, and $d_0=a-1-p-q$. Introduce the excess parameter $v$;
the centered summand and its sum are
\begin{align*}
 r_n(v)&=x_n\frac{\Gamma(n+a)\Gamma(n+b)\Gamma(n+c)
                           \Gamma(n+d_0-v)}
 {\Gamma(n+1)\Gamma(n+1+a-b)\Gamma(n+1+a-c)
                           \Gamma(n+b+c+v)},\\
 G(v)&=\frac{\Gamma(b)\Gamma(c)\Gamma(d_0-v)\Gamma(v)}
 {2\Gamma(d_0)\Gamma(b+v)\Gamma(c+v)}.
\end{align*}
The centered Stirling expansion \cite{DLMFGammaAsymptotic} is
$r_n(v)=x_n^{-1-2v}(1+C_1(v)x_n^{-2}+O(x_n^{-4}))$.
If $\eta=a/2$ and
\[
 (\alpha_1,\alpha_2,\alpha_3,\alpha_4)
 =(\eta,1+p-\eta,1+q-\eta,\eta-1-p-q-v),
\]
then $C_1(v)=-\tfrac13\sum_i B_3(\alpha_i)$.
Subtract the first two centered terms. Normal convergence continues
the resulting sum from $\Re v>0$ to a neighborhood of $v=-1$.
At that point $r_n(-1)=x_nQ_n$ and
\[
 \rho=C_1(-1)=-d_0pq,\qquad
 C_1'(-1)=B_2(\eta-p-q),\qquad
 \zeta(-1,\eta)=-\tfrac12B_2(\eta).
\]
The finite part of $G(v)$ is
\[
 \frac\rho2\{\psi(2)-\psi(a-p-q)-\psi(p)-\psi(q)\}.
\]
The polar Hurwitz counterterm contributes
$C_1'(-1)/2-\rho\psi(\eta)$. Finally,
\[
 \frac12C_1'(-1)+\zeta(-1,\eta)
                  =-\frac12d_0(p+q).
\]
Using $\psi(p)=\psi(1+p)-1/p$ and its $q$ counterpart cancels
these last terms and gives \eqref{bell:master}.
The calculation first applies for $a>1$ and generic small $p,q$.
The bracket is uniformly $O(n^{-3})$ on compact parameter sets, and its
Gamma arguments avoid poles in a neighborhood of every $(a,0,0)$
with $\Re a>0$. Analytic continuation therefore gives the asserted
domain. Cauchy's formula on a slightly larger parameter polydisc
justifies every fixed derivative and includes the loci $p=0$ or $q=0$.
\end{proof}

Let $L_2,L_3,\ldots$ temporarily be independent indeterminates of
weights $2,3,\ldots$. Define
\begin{equation}
 \mathcal B_{r,s}
 =r!s![p^rq^s]\exp\left(
  \sum_{k\ge2}\frac{L_k}{k!}\{(p+q)^k-p^k-q^k\}\right),
 \qquad r,s\ge1.
 \label{bell:definition}
\end{equation}
Only finitely many coefficients enter each polynomial. Taylor expansion
of the logarithm of \eqref{bell:Q} gives
\[
 \log Q_n(p,q;a)
 =\sum_{i,j\ge1}L_{i+j}(n;a)\frac{p^iq^j}{i!j!}.
\]
Thus \eqref{bell:definition} specializes to the mixed derivative of
$Q_n$ at the origin. We retain the distinction between an indeterminate
$L_k$ and its specialization \eqref{bell:lambda} whenever dimensions
are at issue.

\subsection{Formal rank at every weight}

Let $\mathcal P_w$ be the rational vector space spanned by all monomials
of weight $w$ in $L_2,L_3,\ldots$. Let $p(w)$ be the ordinary integer
partition number, with $p(0)=1$. Its dimension is
$p(w)-p(w-1)$, since these are exactly the partitions without a part one.
Set
\[
 \mathcal V_w=\Span_{\Q}\{\mathcal B_{r,w-r}:1\le r\le\lfloor w/2\rfloor\}.
\]
The symmetry exchanging the two Gamma parameters has already been
included: $\mathcal B_{r,s}=\mathcal B_{s,r}$. The elementary linear
row is already present, since $\mathcal B_{1,w-1}=L_w$.

\begin{theorem}[Complete formal rank and finite membership test]
\label{bell:rank}
For every integer $w\ge2$,
\begin{equation}
 \dim_{\Q}\mathcal V_w=\lfloor w/2\rfloor,\qquad
 \operatorname{codim}_{\mathcal P_w}\mathcal V_w
       =p(w)-p(w-1)-\lfloor w/2\rfloor.
 \label{bell:rankformula}
\end{equation}
For a monomial $L_\lambda=\prod_{k\ge2}L_k^{\nu_k}$,
$\sum k\nu_k=w$, its coefficient in row $r$ is exactly
\begin{equation}
 M_w(r,\lambda)
 =\frac{r!(w-r)!}{\prod_k\nu_k!(k!)^{\nu_k}}
 [t^r]\prod_{k\ge2}\{(1+t)^k-1-t^k\}^{\nu_k}.
 \label{bell:matrixformula}
\end{equation}
Consequently rational row reduction of this explicit finite matrix
decides membership in the stated moment space at any prescribed weight.
\end{theorem}

\begin{proof}
Expanding the exponential in \eqref{bell:definition}, and then setting
$q=1$, proves \eqref{bell:matrixformula}. A monomial with $\ell$ factors
requires at least $\ell$ powers of both $p$ and $q$, because every
term in the exponent is divisible by $pq$. Hence a row with first
index $r$ contains no monomial with more than $r$ factors.

For $1\le j\le\lfloor w/2\rfloor$, select the monomial
\[
 U_j=L_2^{j-1}L_{w-2j+2}.
\]
It has exactly $j$ factors, including when it is $L_2^j$. The
coefficient of $U_j$ vanishes in every row $r<j$. In row $r=j$ it is
strictly positive: choose the linear term of each factor
$(1+t)^k-1-t^k$, whose coefficients are positive. The square minor
formed by these columns is triangular with nonzero diagonal. The
rows are therefore independent. There are exactly $\lfloor w/2\rfloor$
of them, which proves both dimension formulas and the membership claim.
\end{proof}

\begin{remark}[The scope of the obstruction]
Nonmembership in $\mathcal V_w$ means that these particular Gamma
derivative rows and their parameter-exchange symmetry do not isolate
the desired polynomial. It does not imply that its sum lacks a closed
form, nor that evaluated special values are arithmetically independent.
Differentiating in $a$, shifting summation indices, or using a different
hypergeometric identity can supply additional relations outside this space.
\end{remark}

\subsection{The requested order-six and order-eight elimination matrices}

At $w=6$ order the columns as
$(L_6,L_2L_4,L_3^2,L_2^3)$ and the rows as $(1,5),(2,4),(3,3)$.
Then
\begin{equation}
 M_6=\begin{pmatrix}1&0&0&0\\1&8&6&0\\1&9&9&6\end{pmatrix},\qquad
 E_6=\begin{pmatrix}1&0&0\\-1/6&1/2&-1/3\\1/18&-1/2&4/9\end{pmatrix},
 \label{bell:sixmatrix}
\end{equation}
and exact multiplication gives
\begin{equation}
 E_6M_6=\begin{pmatrix}1&0&0&0\\0&1&0&-2\\0&0&1&8/3\end{pmatrix}.
 \label{bell:sixrref}
\end{equation}
Thus the combinations $L_2L_4-2L_2^3$ and
$L_3^2+\tfrac83L_2^3$ are evaluated by the family. None of the three
nonlinear monomials is individually in its row space. The column vector
$(0,2,-8/3,1)^T$ annihilates every row and has a nonzero entry in
each nonlinear position, giving a short exact certificate of this assertion.

At $w=8$, take the columns
\[
 (L_8,L_2L_6,L_3L_5,L_4^2,L_2^2L_4,L_2L_3^2,L_2^4)
\]
and the rows $(1,7),(2,6),(3,5),(4,4)$. The matrix is
\begin{equation}
 M_8=\begin{pmatrix}
 1&0&0&0&0&0&0\\
 1&12&30&20&0&0&0\\
 1&15&45&30&60&90&0\\
 1&16&48&34&72&144&24
 \end{pmatrix}.
 \label{bell:eightmatrix}
\end{equation}
An explicit elimination matrix and its result are
\begin{align}
 E_8&=\begin{pmatrix}
 1&0&0&0\\-1/6&1/2&-1/3&0\\
 1/90&-1/6&22/45&-1/3\\1/30&0&-8/15&1/2
 \end{pmatrix},\notag\\
 E_8M_8&=\begin{pmatrix}
 1&0&0&0&0&0&0\\
 0&1&0&0&-20&-30&0\\
 0&0&1&0&16/3&-4&-8\\
 0&0&0&1&4&24&12
 \end{pmatrix}.
 \label{bell:eightrref}
\end{align}
The three free columns give the following basis of the annihilator:
\begin{equation}
 \begin{pmatrix}0\\20\\-16/3\\-4\\1\\0\\0\end{pmatrix},\qquad
 \begin{pmatrix}0\\30\\4\\-24\\0\\1\\0\end{pmatrix},\qquad
 \begin{pmatrix}0\\0\\8\\-12\\0\\0\\1\end{pmatrix}.
 \label{bell:eightann}
\end{equation}
Every nonlinear column has a nonzero coordinate in at least one of
these annihilators. Therefore $L_8$ is the only individual monomial
in this weight that the specified row space isolates. Equations
\eqref{bell:sixmatrix}--\eqref{bell:eightann} answer the requested finite
elimination problem with exact rational certificates.

\subsection{Summation reduces the image to three functions}

For $w\ge6$, write
\begin{equation}
 T_w(a)=(a-1)\psi^{(w-2)}(a)+(w-2)\psi^{(w-3)}(a),
 \qquad \epsilon_w=(-1)^w,
 \label{bell:T}
\end{equation}
and let $S_r(a)$ denote
$\sum_{n\ge0}x_n\mathcal B_{r,w-r}(n;a)$.
No subtraction is needed at these weights. More generally, every
monomial of weight $w\ge4$ has an absolutely convergent weighted sum:
the polygamma series gives
$L_{2j}=O(n^{-2j})$ and $L_{2j+1}=O(n^{-2j})$, locally uniformly in
$a$. The smallest resulting decay exponent is at least four, before
multiplication by $x_n$.

\begin{theorem}[Exact summation image and all its zero rows]
\label{bell:image}
For $w\ge6$ the moment functions are
\begin{align}
 S_1(a)&=\frac{w-1}{2}\left\{
 \epsilon_wT_w(a)+(a-1)\psi^{(w-2)}(1)
                  -(w-2)\psi^{(w-3)}(1)\right\},\label{bell:S1}\\
 S_2(a)&=\epsilon_w(w-2)T_w(a)-(w-2)\psi^{(w-3)}(1),\label{bell:S2}\\
 S_r(a)&=\frac{\epsilon_wr(w-r)}2T_w(a),
 \qquad 3\le r\le\lfloor w/2\rfloor.\label{bell:Sr}
\end{align}
The linear map from $\mathcal V_w\otimes_{\Q}\C$ to holomorphic
functions of $a$ obtained by weighted summation has rank exactly three.
Its kernel has dimension $\lfloor w/2\rfloor-3$, and a rational basis is
\begin{equation}
 \frac{\mathcal B_{r,w-r}}{r(w-r)}
       -\frac{\mathcal B_{3,w-3}}{3(w-3)},
 \qquad 4\le r\le\lfloor w/2\rfloor.
 \label{bell:kernelbasis}
\end{equation}
Every one of these polynomials has identically zero weighted sum for
$\Re a>0$.
\end{theorem}

\begin{proof}
Put $b_{i,j}=\partial_p^i\partial_q^jb_a(0,0)$. For $i+j=k\ge1$,
\[
 b_{i,j}=-\tfrac12\left\{
 \mathbf1_{j=0}\psi^{(i)}(1)+\mathbf1_{i=0}\psi^{(j)}(1)
                     +(-1)^k\psi^{(k)}(a)\right\}.
\]
Differentiating \eqref{bell:master} gives
\[
 S_r(a)=-r(w-r)\{A b_{r-1,w-r-1}
 -(r-1)b_{r-2,w-r-1}-(w-r-1)b_{r-1,w-r-2}\}.
\]
Terms with a negative $b$-index are omitted; their prefactors are zero.
The counterterm has total degree only three, and therefore contributes
nothing here. The cases $r=1$, $r=2$, and $r\ge3$ are precisely
\eqref{bell:S1}--\eqref{bell:Sr}.

To prove exact rank, not merely an upper bound, let $a$ tend to positive
infinity. The standard polygamma asymptotic yields
\[
 T_w(a)=\epsilon_w(w-4)!a^{-(w-3)}+O(a^{-(w-2)}).
\]
The first row has a nonzero linear term, the second has a nonzero
constant limit, and the third is a nonzero function tending to zero.
In any constant linear relation, these three successive behaviors force
the coefficients of $S_1$, $S_2$, and $S_3$ to vanish. The image rank is
exactly three. The displayed kernel rows are independent in the formal
basis of Theorem~\ref{bell:rank}; their number equals the nullity, so
they form the entire kernel. This proves the claim over $\C$ and hence
also its rational version.
\end{proof}

For example, the two evaluable order-six combinations in
\eqref{bell:sixrref} give
\begin{align}
 \sum_{n\ge0}x_n(L_2L_4-2L_2^3)
 &=\frac{T_6(a)}{12}+10(a-1)\zeta(5)-2\zeta(4),\label{bell:sixidentity1}\\
 \sum_{n\ge0}x_n\left(L_3^2+\frac83L_2^3\right)
 &=\frac{5T_6(a)}{36}-\frac{10(a-1)}3\zeta(5)
                                    +\frac{26}3\zeta(4).
 \label{bell:sixidentity2}
\end{align}
Here and below every $L_k$ in a summand means $L_k(n;a)$.

\begin{corollary}[A quartic polygamma identity]
\label{bell:quartic}
For $\Re a>0$,
\begin{equation}
 \sum_{n\ge0}x_n\left\{
 L_8-30L_4^2-120L_2^2L_4-720L_2L_3^2-360L_2^4\right\}=0.
 \label{bell:quarticzero}
\end{equation}
Equivalently,
\begin{align}
 &\sum_{n\ge0}x_n\{L_4^2+4L_2^2L_4+24L_2L_3^2+12L_2^4\}\notag\\
 &\qquad=\frac7{60}\left\{T_8(a)+(a-1)\psi^{(6)}(1)
                                         -6\psi^{(5)}(1)\right\}.
 \label{bell:quarticvalue}
\end{align}
\end{corollary}

\begin{proof}
The polynomial in \eqref{bell:quarticzero} is exactly
$16\mathcal B_{3,5}-15\mathcal B_{4,4}$, by \eqref{bell:eightmatrix}.
Its sum is zero by \eqref{bell:Sr}. Move the linear term to the other
side and apply \eqref{bell:S1}. All the individual sums converge,
so this rearrangement needs no regularized interpretation.
\end{proof}

\subsection{A harmonic-tail and multiple-polylogarithm consequence}

For $k\ge1$, define the centered cubic tail
\[
 C_3(k)=2\{\zeta(3)-H_{k-1}^{(3)}\}-k^{-3}
       =k^{-3}+2\sum_{j>k}j^{-3}.
\]
The next identity is a consequence of the order-eight null row, rather
than an independently asserted new evaluation in the classical Euler-sum
literature.

\begin{corollary}[A reciprocal-weight cubic-tail square]
\label{bell:harmonic}
One has the ordinary convergent identities
\begin{align}
 \sum_{k=1}^{\infty}\frac{C_3(k)^2}{k}&=2\zeta(7),\label{bell:tailsquare}\\
 \Li_{3,4}(1)+2\Li_{3,3,1}(1)+\Li_{6,1}(1)&=\frac14\zeta(7).
 \label{bell:mzv}
\end{align}
The multiple polylogarithms use strictly decreasing positive indices.
\end{corollary}

\begin{proof}
At $a=2$, set $k=n+1$. The polygamma functional equation gives
\[
 x_n=k,\quad L_2=-k^{-2},\quad L_4=-6k^{-4},\quad
 L_8=-5040k^{-8},\quad L_3=2C_3(k).
\]
Substituting in \eqref{bell:quarticzero} yields
$-5760\zeta(7)+2880\sum_k C_3(k)^2/k=0$, proving
\eqref{bell:tailsquare}. Since all terms in the defining tail are
positive, expansion and rearrangement are justified by monotone
convergence (or directly by absolute convergence). In decreasing-index
notation that expansion is
\[
 \sum_k\frac{C_3(k)^2}{k}
 =\zeta(7)+4\zeta(3,4)+8\zeta(3,3,1)+4\zeta(6,1).
\]
Combining the two expressions proves \eqref{bell:mzv}.
\end{proof}

The rank theorem explains both the usefulness and the limitations of
the two-parameter family. It supplies many exact nonlinear combinations
and an increasing number of zero-sum relations, while the explicit
annihilators at weights six and eight explain why an individual
unseparated moment still requires additional information.

% END sections/03_bell_moments.tex

% BEGIN sections/04_harmonic_powers.tex
\section{A finite Euler-sum space for every centered harmonic moment}
\label{hpnew:section}

The recent raw-power report proves that the centered series
\[
 F_p(s,a)=\sum_{n=0}^{\infty}\frac{H_n^p}{(n+a)^s},\qquad
 H_n=\sum_{j=1}^n\frac1j,\quad H_0=0,
\]
is regular at $s=0,-2,-4,\ldots$ when $a=1/2$.
Its Question 10 asks for a generator that avoids listing every
composition of $p$, with the complete sequence $F_p(-2,1/2)$ as a
first target \cite{RepoExactResonant}. We give such a generator for
\emph{every} nonnegative even moment. More strongly, for fixed $p$ all
these values lie in one finite rational span of convergent Euler sums,
independent of the moment order. These are exact analytic identities;
no assertion of arithmetic independence or literature-wide priority is
made. They concern values at regular spectral points, not derivatives
of a formula known only at those points.

\subsection{An exponent parameter with the required joint regularity}

Introduce the auxiliary exponent family
\begin{equation}\label{hpnew:exponent}
 \mathscr F(s,u)=\sum_{n=0}^{\infty}
       \frac{\exp(uH_n)}{(n+1/2)^s},\qquad
                    \Re(s-u)>1.
\end{equation}
Positive bases have their real logarithms. This normally convergent
series is jointly holomorphic in its initial domain. At $u=0$ its
$p$th derivative is $F_p(s,1/2)$ for $p\ge1$; the zeroth derivative is
$\zeta(s,1/2)$. Let $B_j(x)$ denote the Bernoulli polynomials, and put
\begin{equation}\label{hpnew:coefficients}
 b_j=-\frac{B_{2j}(1/2)}{2j},\qquad
 c_j(u)=[z^j]\exp\left(u\sum_{\ell\ge1}b_\ell z^\ell\right).
\end{equation}
The latter definition is formal: each $c_j$ is a finite polynomial.
No convergence of the infinite Bernoulli series is required.

\begin{lemma}[Joint continuation and the centered divisor pattern]
\label{hpnew:joint}
The family $\mathscr F$ has a joint meromorphic continuation to
$\mathbb C^2$, with possible polar hyperplanes
\begin{equation}\label{hpnew:divisors}
                         s-u=1-2j,\qquad j=0,1,2,\ldots.
\end{equation}
In particular it is jointly holomorphic near every $(-2r,0)$,
$r\ge0$. The function
\begin{equation}\label{hpnew:V}
 V_r(u)=\mathscr F(-2r,u)
       =\sum_{p\ge1}F_p(-2r,1/2)\frac{u^p}{p!}
\end{equation}
is holomorphic for $|u|<1$.
\end{lemma}

\begin{proof}
With $x=n+1/2$, Euler--Maclaurin gives, for every fixed $J$,
\[
 H_n=\log x+\gamma+\sum_{j=1}^Jb_jx^{-2j}
                              +O(x^{-2J-2}).
\]
This is the centered shifted-digamma expansion used in
\cite{RepoExactResonant}. Exponentiating a finite truncation gives
\begin{equation}\label{hpnew:asymptotic}
 \exp(uH_n)=e^{\gamma u}x^u
      \left(\sum_{j=0}^Jc_j(u)x^{-2j}
                         +O_K(x^{-2J-2})\right),
\end{equation}
uniformly for $u$ in a compact set $K$. Fixed $u$ derivatives have
analogous estimates with powers of $\log x$.
Subtract these finitely many terms in \eqref{hpnew:exponent}, starting
at any fixed positive index. The remainder is normally convergent
when $\Re(s-u)>-2J-1$, and the subtracted terms are
$e^{\gamma u}c_j(u)\zeta(s-u+2j,N+1/2)$.
The sole pole of each Hurwitz function is at its first argument $1$
\cite{DLMFHurwitz}. Arbitrarily large $J$ prove the stated joint
continuation. None of the hyperplanes \eqref{hpnew:divisors} passes
through $(-2r,0)$. Joint holomorphy, rather than a formal interchange
at a singular point, therefore identifies all the coefficients in
\eqref{hpnew:V}. Its constant coefficient is
$\zeta(-2r,1/2)=0$.
\end{proof}

\subsection{A triangular generator from ordinary telescoping}

For $k\ge2$ define the normally convergent functions
\begin{align}
 E_k(u)&=\sum_{n\ge1}\frac{\exp(uH_{n-1})}{n^k},
                                  &&\Re u<k-1,\label{hpnew:Ek}\\
 B(u)&=1+\sum_{k\ge2}\frac{u^k}{k!}E_k(u),
                                  &&\Re u<1.\label{hpnew:B}
\end{align}
The letter $B$ without an index in \eqref{hpnew:B} is an auxiliary
function, not a Bernoulli polynomial. Its coefficients involve the
ordinary convergent Euler sums
\begin{equation}\label{hpnew:Euler}
 E_{q,k}=\sum_{n\ge1}\frac{H_{n-1}^q}{n^k},\qquad
 q\ge0,\quad k\ge2;\qquad
 E_k(u)=\sum_{q\ge0}E_{q,k}\frac{u^q}{q!}.
\end{equation}
The Taylor identity in \eqref{hpnew:Euler} is used near $u=0$.
Here $H_{n-1}^0=1$, including at $n=1$.

\begin{theorem}[A finite triangular generator at every centered moment]
\label{hpnew:generator}
For $m=0,1,\ldots$, define $G_m(u)$ recursively by
\begin{align}
 (u+m+1)G_m(u)
 ={}&\sum_{j=0}^{m-1}
 \left\{(-1)^{m+1-j}\binom{m+1}{j}
             -\frac{u^{m+1-j}}{(m+1-j)!}\right\}G_j(u)\notag\\
 &+\frac{u^{m+1}}{(m+2)!}B(u)
 -\sum_{k\ge2}\frac{u^{k+m+1}}{(k+m+1)!}E_k(u).
 \label{hpnew:Grecurrence}
\end{align}
The first sum is empty at $m=0$. All functions are holomorphic on
$|u|<1$, and
\begin{equation}\label{hpnew:Vrecurrence}
 \boxed{\displaystyle
 V_r(u)=\sum_{j=0}^{2r}\binom{2r}{j}
                     \left(-\frac12\right)^{2r-j}G_j(u).}
\end{equation}
Thus $2r+1$ triangular steps determine the exponential generating
function for all $F_p(-2r,1/2)$. For a specified $p$, only finitely
many coefficients of finitely many $E_k$ are needed.
\end{theorem}

\begin{proof}
Put $g_n(u)=\exp(uH_{n-1})$. It obeys the exact discrete equation
\[
                  g_{n+1}(u)=e^{u/n}g_n(u).
\]
For $\Re u<0$, telescoping and then an absolutely convergent
exponential expansion give
\[
 -1=\sum_{n\ge1}(g_{n+1}-g_n)
   =u\sum_{n\ge1}\frac{g_n}{n}
                      +\sum_{k\ge2}\frac{u^k}{k!}E_k(u).
\]
Consequently the continued order-one sum is exactly
\begin{equation}\label{hpnew:Eone}
                \sum_{n\ge1}\frac{g_n(u)}n=-\frac{B(u)}u.
\end{equation}
This equality initially concerns ordinary sums in $\Re u<0$.

For a fixed $m\ge0$, work first in $\Re u<-m-1$, where all the
following sums converge absolutely. Multiplying the discrete equation
by $n^{m+1}$, summing to a finite cutoff $N$, and only then
shifting an index yields the following formulas after $N\to\infty$
\begin{align*}
 \sum_{n\ge1}n^{m+1}(g_{n+1}-g_n)
 &=\sum_{j=0}^{m}(-1)^{m+1-j}\binom{m+1}{j}
                                  \sum_{n\ge1}n^jg_n,\\
 \sum_{n\ge1}n^{m+1}(g_{n+1}-g_n)
 &=\sum_{k\ge1}\frac{u^k}{k!}
                                  \sum_{n\ge1}\frac{g_n}{n^{k-m-1}}.
\end{align*}
The upper-end term is $N^{m+1}g_{N+1}=O(N^{m+1+\Re u})$ and
vanishes in the stated half-plane. There is no lower-end boundary
term because the polynomial being shifted has positive degree.
This finite-cutoff argument never separates a difference into two
possibly divergent infinite sums. In the second line separate $k=1$,
$2\le k\le m+1$, $k=m+2$, and $k\ge m+3$.
Using \eqref{hpnew:Eone} for the order-one term and moving the two
copies of $\sum n^mg_n$ to the left gives
\eqref{hpnew:Grecurrence}. In particular its $G_m$ is the
meromorphic continuation of the ordinary sum
$\sum_{n\ge1}n^m\exp(uH_{n-1})$ from $\Re u<-m-1$.

Here are the convergence details for the continuation. On a compact
subset of $\Re u<1$, choose $U<1$ with $\Re u\le U$.
The two-sided estimate $H_{n-1}=\log n+O(1)$ gives
$|g_n(u)|\le Cn^U$, uniformly, including when $\Re u$ is negative. Hence all $E_k(u)$, $k\ge2$, are
bounded by a common summable majorant, and the factorial denominators
make both infinite $k$ sums normally convergent. Fixed derivatives in
$u$ are justified on slightly larger compact subsets, or directly by
inserting powers of $H_{n-1}$. The recurrence therefore continues
$G_m$ meromorphically to $\Re u<1$, with possible poles only at
$-1,-2,\ldots,-m-1$. In particular it is holomorphic on $|u|<1$.

Finally, for $\Re u<-2r-1$, an ordinary change of index in the
absolutely convergent defining series gives
\[
 \mathscr F(-2r,u)=\sum_{n\ge1}(n-1/2)^{2r}g_n(u)
   =\sum_{j=0}^{2r}\binom{2r}{j}(-1/2)^{2r-j}G_j(u).
\]
Meromorphic continuation in $u$ proves \eqref{hpnew:Vrecurrence};
Lemma~\ref{hpnew:joint} identifies its Taylor coefficients with the
claimed regular spectral values.
\end{proof}

\begin{remark}[Why the auxiliary values cannot be specialized in the other order]
\label{hpnew:orderwarning}
Induction in \eqref{hpnew:Grecurrence} gives $G_m(0)=0$ for every
$m$. For example $G_1(0)=0$, whereas $\zeta(-1)=-1/12$.
This is consistent: $G_m$ first fixes the spectral argument at $-m$
and then continues the exponent $u$. The uncentered family has polar hyperplanes $s-u=1-j$.
Even a single such divisor can prevent joint holomorphy when its
numerator vanishes at the specialization point. For example its
local subtraction includes
\[
 e^{\gamma u}\left(\frac{u^2}{8}-\frac{u}{12}\right)
                                      \zeta(s-u+2).
\]
At $s=-1$ and then $u\to0$, this contributes $1/12$; at $u=0$
first, it contributes zero. It is therefore invalid to replace
$G_m(0)$ by the separately specialized zeta value.
The centered combination \eqref{hpnew:Vrecurrence} is justified by
joint holomorphy, which was proved before coefficient extraction.
\end{remark}

\subsection{A finite arithmetic space independent of the moment order}

\begin{theorem}[Uniform finite Euler-sum span]
\label{hpnew:uniformspan}
For $p\ge1$, set
\begin{equation}\label{hpnew:space}
 \mathcal E_{p-1}=\operatorname{span}_{\mathbb Q}
 \left(\{1\}\cup
   \{E_{q,k}:q\ge0,\ k\ge2,\ q+k\le p-1\}\right).
\end{equation}
Then
\begin{equation}\label{hpnew:spanconclusion}
 \boxed{\displaystyle
             F_p(-2r,1/2)\in\mathcal E_{p-1}
                    \quad\hbox{for every }r\ge0.}
\end{equation}
Every value has a finite multiple-zeta expression using weights at
most $p-1$, with no Euler constant required. In particular every
centered even-negative moment of $H_n$ or $H_n^2$ is rational.
For $p\ge2$, the displayed spanning list has at most
$1+\binom{p-1}{2}$ elements; this is a spanning bound, not a claim
that the periods in the list are independent.
\end{theorem}

\begin{proof}
The recurrence is linear in the functions $E_k(u)$, and all rational
coefficients being expanded are regular at $u=0$.
An occurrence of $E_k$ first has a factor $u^{k+m+1}$, either in
the last sum of \eqref{hpnew:Grecurrence} or through $B$ in its
preceding term. Its exponent is at least $k+1$.
Induction on $m$ therefore shows that the coefficient of $u^p$ in
any $G_m$, and hence in any $V_r$, is a rational linear combination
of $1$ and $E_{q,k}$ with $q+k\le p-1$.

For completeness, the finite multiple-zeta reduction is
\begin{equation}\label{hpnew:MZV}
 E_{q,k}=\sum_{\substack{(r_1,\ldots,r_d)\text{ a composition of }q}}
       \frac{q!}{r_1!\cdots r_d!}\,
                          \zeta(k,r_1,\ldots,r_d).
\end{equation}
At $q=0$ the empty composition gives $E_{0,k}=\zeta(k)$.
To prove the formula, expand $H_{n-1}^q$, order its distinct indices,
and group equal indices into blocks of sizes $r_1,\ldots,r_d$.
The displayed multinomial is their multiplicity. Each resulting MZV
converges because its first index is $k\ge2$, and its weight is
$q+k$. Formula \eqref{hpnew:MZV} is an optional final conversion:
it is not necessary for computing the generator itself.
\end{proof}

This uniform space is smaller information than an unspecified
multiple-zeta algebra at each separate moment. Increasing $r$ changes
rational coefficients in one fixed finite list; it does not require
Euler sums of increasing weight.

\subsection{Two all-moment Bernoulli identities}

\begin{corollary}[The first two powers at every centered even moment]
\label{hpnew:Bernoulli}
For every $r\ge0$,
\begin{align}
 F_1(-2r,1/2)&=\frac12 B_{2r}(1/2),\label{hpnew:firstpower}\\
 F_2(-2r,1/2)&=-\sum_{j=0}^{r}\binom{2r}{2j}
       \frac{B_{2j}B_{2r-2j}(1/2)}{2j+1}.
                         \label{hpnew:secondpower}
\end{align}
\end{corollary}

\begin{proof}
The classical generating functions needed here are particularly short.
Write $L(t)=-\log(1-e^{-t})$, and define
$A_p(t)=(1-e^{-t})\sum_{n\ge0}H_n^pe^{-nt}$.
Then $A_1=L$ and $A_2=L^2+\Li_2(e^{-t})$.
Consequently
\[
                         A_2'(t)=-\coth(t/2)L(t).
\]
The odd powers in the local expansion of $L$ consist only of $t/2$:
\[
 L(t)=-\log t+\frac t2
                  -\log\left(\frac{\sinh(t/2)}{t/2}\right).
\]
Here ``odd powers'' refers to the power--log expansion at $t=0$,
with no branch continuation to negative $t$ being assumed.
It follows that the odd-power part of $A_2$ is
\[
 -\int_0^t\frac v2\coth(v/2)\,dv
       =-\sum_{j\ge0}\frac{B_{2j}}{(2j)!(2j+1)}t^{2j+1}.
\]
The centered Mellin kernel is $A_p(t)/(2\sinh(t/2))$.
Its log-free coefficient of $t^{2r}$, multiplied by $(2r)!$, is
$F_p(-2r,1/2)$. Indeed that coefficient gives the simple Mellin
pole at $s=-2r$, and the simple zero of $1/\Gamma(s)$ cancels
it; all holomorphic Mellin terms vanish at the specialization.
Regularity proved in Lemma~\ref{hpnew:joint} eliminates higher
logarithmic poles. Finally use
\[
 \frac1{2\sinh(t/2)}=
             \sum_{j\ge0}\frac{B_{2j}(1/2)}{(2j)!}t^{2j-1}.
\]
Multiplying the two convergent local series gives
\eqref{hpnew:firstpower} and \eqref{hpnew:secondpower}.
\end{proof}

\subsection{A sufficient parity class for generalized harmonic products}

There is also a direct extension toward Question 11 of
\cite{RepoExactResonant}. Define
$H_n^{(q)}=\sum_{k=1}^n k^{-q}$ and
$U_q(s)=\sum_{n\ge0}H_n^{(q)}(n+1/2)^{-s}$.

\begin{corollary}[Odd harmonic orders and finite ordinary-zeta values]
\label{hpnew:oddorders}
Let $P$ be any polynomial in finitely many variables. The Dirichlet
series with numerator
\[
                 P(H_n,H_n^{(3)},H_n^{(5)},\ldots,H_n^{(2d+1)})
\]
is holomorphic at every $s=-2r$, $r\ge0$. This is a sufficient
class of numerator polynomials, not a necessity classification.
For a single positive odd integer $q$, the exact values are
\begin{equation}\label{hpnew:oddformula}
 \boxed{\displaystyle
 U_q(-2r)=-\frac1{2r+1}\sum_{j=0}^{r}\binom{2r+1}{2j}
          B_{2j}(1/2)\,
                         \zeta(q-2r-1+2j).}
\end{equation}
For fixed odd $q\ge3$, all these moments belong to the same space
$\operatorname{span}_{\mathbb Q}\{1,\zeta(2),\zeta(4),\ldots,
\zeta(q-1)\}$. At $q=1$, formula \eqref{hpnew:oddformula} reduces
to \eqref{hpnew:firstpower}.
\end{corollary}

\begin{proof}
For $q\ge3$ and $x=n+1/2$, use
$H_n^{(q)}=\zeta(q)-\zeta(q,x+1/2)$ and the midpoint expansion
\[
 \zeta(q,x+1/2)=\frac{x^{1-q}}{q-1}
 +\sum_{j=1}^{J}\frac{B_{2j}(1/2)(q)_{2j-1}}{(2j)!}
                                       x^{1-q-2j}
 +O(x^{-q-2J-1}).
\]
When $q$ is odd, every displayed inverse power is even. Together
with the even inverse-power expansion of $H_n$ following its
logarithm, this shows that any stated polynomial has an asymptotic
expansion containing only $x^{-2j}(\log x)^h$. Finite subtraction
in its Dirichlet series leaves possible poles only at $s=1-2j$.
This proves regularity at every nonpositive even integer.

To prove the value formula, first take $\Re q>2r+2$. The identity
$U_q(s)-\zeta(q)\zeta(s,1/2)
=-\sum_{n\ge0}\zeta(q,n+1)(n+1/2)^{-s}$ continues the difference
to $s=-2r$ by an absolutely convergent tail sum. Since
$\zeta(-2r,1/2)=0$, ordinary summation reversal gives
\[
 U_q(-2r)=-\sum_{k\ge1}k^{-q}
                         \sum_{n=0}^{k-1}(n+1/2)^{2r}.
\]
The centered Faulhaber identity is
\[
 \sum_{n=0}^{k-1}(n+1/2)^{2r}
 =\frac{B_{2r+1}(k+1/2)}{2r+1}
 =\frac1{2r+1}\sum_{j=0}^{r}\binom{2r+1}{2j}
                              B_{2j}(1/2)k^{2r+1-2j}.
\]
Insertion proves \eqref{hpnew:oddformula} in this half-plane.

The same tail expansion supplies joint meromorphic continuation
in $(s,q)$; its possible moving divisors are $s+q=2-2j$.
None passes through $(-2r,q)$ for an odd positive integer $q$.
The apparent pole at $q=1$ cancels: near $s=-2r$ the combination
\[
 \zeta(q)\zeta(s,1/2)-\frac{\zeta(s+q-1,1/2)}{q-1}
\]
is jointly holomorphic after writing it as
$(\zeta(q)-(q-1)^{-1})\zeta(s,1/2)
-(\zeta(s+q-1,1/2)-\zeta(s,1/2))/(q-1)$.
The remaining finite subtractions and the normal remainder have
no divisor there. Thus meromorphic continuation of the value formula
really gives the separately specified odd-order functions, including
$H_n^{(1)}=H_n$. Finally, every zeta argument in
\eqref{hpnew:oddformula} is even; negative even arguments give zero
and $\zeta(0)=-1/2$, proving the claimed fixed span.
\end{proof}

For example,
\begin{align}\label{hpnew:thirdorder}
 U_3(0)&=-\zeta(2),&
 U_3(-2)&=\frac16+\frac1{12}\zeta(2),&
 U_3(-4)&=-\frac1{12}-\frac7{240}\zeta(2),\\
 U_5(0)&=-\zeta(4),&
 U_5(-2)&=-\frac13\zeta(2)+\frac1{12}\zeta(4),\notag\\
 U_5(-4)&=\frac1{10}+\frac16\zeta(2)-\frac7{240}\zeta(4).
                                                               \notag
\end{align}
The odd-order restriction should be retained. At an even-order
specialization a vanishing coefficient can still meet a moving divisor,
so continuing the fixed spectral value in $q$ need not agree with
specializing $q$ before taking the spectral value. No such ambiguity
occurs at the odd orders proved here.

\subsection{Compact generators at zero, minus two, and minus four}

Put $Q(u)=V_0(u)$, $V(u)=V_1(u)$, and $W(u)=V_2(u)$.
The first three moment rows have the following compact forms.
\begin{corollary}[Three complete rows in the harmonic power]
\label{hpnew:threerows}
For $|u|<1$,
\begin{align}
 (u+1)Q(u)
 &=\frac u2+\frac12\sum_{k\ge2}(k-1)
                  \frac{u^{k+1}}{(k+1)!}E_k(u),\label{hpnew:Q}\\
 (u+3)V(u)
 &=\frac{u^3-3u-3}{12}Q(u)-\frac{u^3}{24}\notag\\
 &\hspace{4mm}-\frac1{24}\sum_{k\ge2}(k-1)(k+1)(k+6)
                  \frac{u^{k+3}}{(k+3)!}E_k(u),\label{hpnew:minus2}\\
 (u+5)W(u)
 &=-\frac{u^3+18u^2+90u+150}{60}V(u)\notag\\
 &\hspace{4mm}-\frac{u^3+6u^2+15u+15}{240}Q(u)\notag\\
 &\hspace{4mm}+\frac1{120}\sum_{k\ge2}k(k-1)(k+1)
                  \frac{u^{k+5}}{(k+5)!}E_k(u).
                                      \label{hpnew:minus4}
\end{align}
\end{corollary}

\begin{proof}
The $m=0$ case of \eqref{hpnew:Grecurrence}, followed by substitution
of \eqref{hpnew:B}, is \eqref{hpnew:Q}. For the next row use
\[
 V=G_2-G_1+\frac14G_0.
\]
The $G_1$ coefficient is $-u(u+2)/2$ before its substitution, so its
apparent denominator $u+2$ cancels. Substituting $B$ gives the factor
\[
 \frac{k+1}{2(k+3)!}-\frac1{24k!}
       =-\frac{(k-1)(k+1)(k+6)}{24(k+3)!},
\]
which proves \eqref{hpnew:minus2}.
For the last row substitute $m\le4$ into
\[
 W=G_4-2G_3+\frac32G_2-\frac12G_1+\frac1{16}G_0.
\]
Collect $Q$, $V$, the constant part of $B$, and the coefficient of
each $E_k$ separately. The first three are respectively
$-(u^3+6u^2+15u+15)/(240(u+5))$,
$-(u^3+18u^2+90u+150)/(60(u+5))$, and zero; the last is
$k(k-1)(k+1)u^{k+5}/(120(u+5)(k+5)!)$.
These rational simplifications prove \eqref{hpnew:minus4}. The
verification script replays them with symbolic $u$ and $k$.
\end{proof}

The zero row \eqref{hpnew:Q} recovers the recurrence already proved
in \cite{RepoExactResonant}. Equations \eqref{hpnew:minus2} and
\eqref{hpnew:minus4} determine two further complete rows, while
Theorem~\ref{hpnew:generator} supplies every remaining row.
For an explicit finite coefficient version, write
$Q_p=F_p(0,1/2)$, $V_p=F_p(-2,1/2)$ and $Q_0=V_0=0$.
For $p\ge1$, \eqref{hpnew:minus2} says
\begin{align}\label{hpnew:finiteV}
 3V_p+pV_{p-1}
 ={}&\frac{p(p-1)(p-2)}{12}Q_{p-3}
       -\frac p4Q_{p-1}-\frac14Q_p-\frac14\mathbf1_{p=3}\notag\\
 &-\frac1{24}\sum_{k=2}^{p-3}(k-1)(k+1)(k+6)
                     \binom p{k+3}E_{p-k-3,k}.
\end{align}
The first term is omitted for $p<3$, and an empty sum is zero.
There is no infinite coefficient calculation at a fixed $p$.

For example, the next fifth-power identity is
\begin{equation}\label{hpnew:fifth}
 \boxed{\displaystyle
 F_5(-2,1/2)=-\frac{200}{81}-\frac{23}{54}\zeta(2)
                  +\frac5{12}\zeta(3)-\frac{11}{8}\zeta(4).}
\end{equation}
At the next spectral point the first five powers are
\begin{align}
 F_1(-4,1/2)&=\frac7{480},&
 F_2(-4,1/2)&=\frac{19}{3600},\notag\\
 F_3(-4,1/2)&=\frac{493}{18000}+\frac7{480}\zeta(2),
                                                      \label{hpnew:table4}\\
 F_4(-4,1/2)&=-\frac{6479}{67500}+\frac{19}{1800}\zeta(2)
                              +\frac7{80}\zeta(3),\notag\\
 F_5(-4,1/2)&=\frac{98437}{202500}+\frac{643}{5400}\zeta(2)
                +\frac{19}{240}\zeta(3)+\frac{77}{160}\zeta(4).
                                                        \notag
\end{align}
Coefficient extraction from \eqref{hpnew:Q}--\eqref{hpnew:minus4}
proves these evaluations. Beyond $E_{0,k}=\zeta(k)$, the only
low-weight reductions needed are the classical identities
$E_{1,2}=\zeta(3)$, $E_{1,3}=\zeta(4)/4$ and
$E_{2,2}=11\zeta(4)/4$, also used in \cite{RepoExactResonant}.

\subsection{Exact exponent singularities and reproducible checks}

The apparent negative even denominators in the triangular recursion
cancel in every centered row. This has a stronger analytic explanation.
\begin{corollary}[Exact negative poles of the exponent generator]
\label{hpnew:upoles}
On $\Re u<1$, the function $V_r$ has exactly the simple poles
\[
                     u=-2r-1,-2r+1,\ldots,-1.
\]
At $u_j=2j-2r-1$, $0\le j\le r$, its residue is
\begin{equation}\label{hpnew:residue}
             \operatorname{Res}_{u=u_j}V_r(u)
                        =-e^{\gamma u_j}c_j(u_j).
\end{equation}
Its Taylor series at zero therefore has exact radius $1$.
\end{corollary}

\begin{proof}
The finite subtraction in \eqref{hpnew:asymptotic} leaves precisely
these possible poles in the indicated half-plane. The Hurwitz pole is
$\zeta(1-(u-u_j),N+1/2)$, with residue $-1$ in $u$, proving
\eqref{hpnew:residue}. To check that none is removable, use
$B_{2j}(1/2)=(2^{1-2j}-1)B_{2j}$. Thus
$b_j=(-1)^{j+1}|b_j|$. For a negative real $u$, every term in the
coefficient $c_j(u)$ has sign $(-1)^j$, and at least one term is
nonzero; also $c_0=1$. Hence every displayed residue is nonzero.
The pole at $-1$ fixes the exact Taylor radius.
\end{proof}

The accompanying \src{code/verify_harmonic_powers.py} performs exact
rational checks of the three generator rows, extracts the displayed
values, checks the first two Bernoulli formulas and the centered
Faulhaber identity through multiple moment orders, and evaluates six
odd generalized-harmonic checkpoints. Its independent numerical
calculation starts from the
original harmonic-power sums, subtracts their centered Euler--Maclaurin
expansion, and evaluates the remaining Hurwitz order derivatives.
Two different truncation lengths are compared. This diagnostic tests
signs and constants without reusing the triangular generator; it is
not a proof of an infinite identity or a certified numerical error bound.
The proof is the normally convergent telescoping argument and the
joint continuation established above. The first regular spectral
coefficient beyond these centered values still carries global Mellin
information; the theorem makes no unsupported claim to eliminate it.

% END sections/04_harmonic_powers.tex

% BEGIN sections/05_collisions.tex
\section{An exact geometric collision formula and harmonic hierarchy laws}
\label{hc:section}

The incoming mixed-spectral report already determines the collision of
three undifferentiated Stieltjes factors at shifts $(0,a\epsilon,b\epsilon)$
for fixed $0<a<b$. It also supplies the complete spectral subtraction
at a fixed, arbitrarily multiple collision \cite{RepoMixed}. These are
substantial antecedents. A moving configuration such as
$(0,\epsilon,\epsilon+\epsilon^2)$ is a different question: inserting
$\epsilon$-dependent rates in a fixed-rate error estimate is not justified
without uniform control. This section proves an exact formula before
any rate specialization. It permits arbitrarily many polygamma factors,
arbitrary argument-derivative orders, and analytic collision arcs with
several different vanishing orders. Its correction is a finite expression
in polygamma values, derivatives, rational functions, and logarithms.
The corrected remainder is holomorphic in all the shifts.

\subsection{The coordinate convention and the local telescoping identity}

Write
\[
 f_r(z)=-\psi^{(r)}(z),\qquad
 n_i=r_i+1,\qquad \kappa_i=(-1)^{r_i}r_i!,\qquad
 h_i(z)=f_{r_i}(1+z).
\]
Thus $f_0=\gamma_0$ and $f_r=\gamma_0^{(r)}$; this section concerns
Stieltjes index zero with arbitrary argument derivatives. In particular,
no claim about all positive Stieltjes indices is hidden in the notation.
The polygamma recurrence \cite{DLMFPolygamma} gives the exact local formula
\begin{equation}\label{hc:one-sided}
 f_{r_i}(\{x\})=h_i(x)+
       \kappa_i\mathbf1_{x>0}x^{-n_i},\qquad 0<|x|<\delta<1.
\end{equation}
The functions $h_i$ are holomorphic for $|x|<1$.

For real shifts $a_1<\cdots<a_d$ sufficiently close to zero, define
\begin{equation}\label{hc:Cdef}
 \mathcal C_{\boldsymbol r}(\boldsymbol a)
 =\FP_x\int_0^1\prod_{i=1}^d f_{r_i}(\{x+a_i\})\,\dd x,
 \qquad
 Q_{\boldsymbol r}=\FP_x\int_0^1\prod_{i=1}^d f_{r_i}(x)\,\dd x.
\end{equation}
At the seam $x=-a_i$ on the circle, the cutoff coordinate is exactly
$x+a_i$, taken positive, and the deleted right interval has ambient
length $\eta_i$. The finite part is the coefficient independent of all
$\eta_i$ and $\log\eta_i$. The cutoffs may be taken independently.
These are pointwise products away from finitely many distinct seams;
there is no multiplication of distributions with overlapping singular
supports. The merged quantity uses the same ambient coordinate $x$.

Put, for $1\le m\le d$,
\begin{equation}\label{hc:Km}
 K_m(z,\boldsymbol a)=
 \frac{\kappa_m}{(z+a_m)^{n_m}}
 \prod_{k<m}h_k(z+a_k)
 \prod_{k>m}f_{r_k}(z+a_k).
\end{equation}
In a fixed neighborhood of the merged seam one has
\begin{equation}\label{hc:telescope}
 \prod_{i=1}^d f_{r_i}(\{x+a_i\})
 =\prod_{i=1}^d h_i(x+a_i)
      +\sum_{m=1}^d\mathbf1_{x>-a_m}K_m(x,\boldsymbol a).
\end{equation}
Indeed, expand the product according to the first factor for which the
singular part in \eqref{hc:one-sided} is selected. If this first index is
$m$ and $x>-a_m$, then $x>-a_k$ for every $k>m$, so all the subsequent
factors may be replaced by their meromorphic germs $f_{r_k}(x+a_k)$.
This telescoping form replaces a sum over $2^d$ subsets by $d$ terms.

For $m<i$ and $1\le p\le n_i$, let $B_{mi,p}(\boldsymbol a)$ be the
coefficient of $(z+a_i)^{-p}$ in the principal part of $K_m$ at $-a_i$.
It is explicitly
\begin{align}
 B_{mi,p}(\boldsymbol a)
 =\frac{\kappa_i}{(n_i-p)!}
 \left.\partial_z^{\,n_i-p}\left\{
 \frac{\kappa_m}{(z+a_m)^{n_m}}
 \prod_{k<m}h_k(z+a_k)
 \prod_{\substack{k>m\\k\ne i}}f_{r_k}(z+a_k)
 \right\}\right|_{z=-a_i}.
 \label{hc:B}
\end{align}
Only finitely many derivatives are required. All arguments after
this substitution are differences of shifts.

\subsection{Exact subtraction, with a jointly holomorphic remainder}

\begin{theorem}[Geometric collision formula]\label{hc:main}
With the convention of \eqref{hc:Cdef}, set
\begin{equation}\label{hc:counter}
 \mathcal A_{\boldsymbol r}(\boldsymbol a)=
 \sum_{m<i}\left\{
 -B_{mi,1}(\boldsymbol a)\log(a_i-a_m)
 +\sum_{p=2}^{n_i}
  \frac{B_{mi,p}(\boldsymbol a)}{(p-1)(a_i-a_m)^{p-1}}
 \right\}.
\end{equation}
There is a function $\mathcal R_{\boldsymbol r}$ holomorphic in all
$a_i$ on a complex neighborhood of $\boldsymbol a=0$ such that, in
the stated real ordering chamber,
\begin{equation}\label{hc:exact}
 \boxed{\quad
 \mathcal C_{\boldsymbol r}(\boldsymbol a)
 =\mathcal A_{\boldsymbol r}(\boldsymbol a)
       +\mathcal R_{\boldsymbol r}(\boldsymbol a),\qquad
 \mathcal R_{\boldsymbol r}(\boldsymbol0)=Q_{\boldsymbol r}.
 \quad}
\end{equation}
The theorem applies to every number of factors and every tuple of
nonnegative derivative orders. It supplies the complete collision
correction, including its finite terms, rather than just its leading
divergences.
\end{theorem}

\begin{proof}
Choose a fixed small $\delta>0$ and represent the circle near its seam
by $x\in(-\delta,\delta)$. The integral outside this neighborhood and
the integral of the all-$h_i$ term in \eqref{hc:telescope} are
holomorphic in the shifts. Consider one $K_m$ term, integrated from
$-a_m$ to $\delta$ with its right endpoint cutoff at $-a_m$.

It is important not to take the limits of its individual principal-part
coefficients, which can diverge when poles coalesce. Define
\[
 D_m(z,\boldsymbol a)=\prod_{i=m}^d(z+a_i)^{n_i},\qquad
 H_m(z,\boldsymbol a)=D_m(z,\boldsymbol a)K_m(z,\boldsymbol a).
\]
The numerator is jointly holomorphic: explicitly it is
\[
 H_m=\kappa_m\prod_{k<m}h_k(z+a_k)
       \prod_{k>m}\bigl\{\kappa_k+(z+a_k)^{n_k}h_k(z+a_k)\bigr\}.
\]
For distinct shifts, write the polar decomposition
\[
 K_m(z,\boldsymbol a)
   =\sum_{i\ge m}\sum_{p=1}^{n_i}
          \frac{B_{mi,p}(\boldsymbol a)}{(z+a_i)^p}
       +W_m(z,\boldsymbol a),
\]
where the notation $B_{mm,p}$ is used here only for the analogous
principal coefficients at the starting point.
The remainder extends jointly holomorphically as all the shifts merge.
For example, on a fixed circle $|w|=\rho>\delta$ inside the domain of
all $H_m$, with the shifts sufficiently small, the residue theorem gives
\begin{equation}\label{hc:Hermite}
 W_m(z,\boldsymbol a)=\frac1{2\pi\ii}
 \int_{|w|=\rho}\frac{H_m(w,\boldsymbol a)}
  {D_m(w,\boldsymbol a)(w-z)}\,\dd w,\qquad |z|<\rho.
\end{equation}
The denominator stays away from zero on this contour, proving the
claimed joint holomorphy. This is the analytic cancellation usually
encoded by confluent Hermite divided differences.

A primitive of the polar part is
\[
 U_m(z,\boldsymbol a)=
 \sum_{i\ge m} B_{mi,1}\log(z+a_i)
 -\sum_{i\ge m}\sum_{p=2}^{n_i}
       \frac{B_{mi,p}}{p-1}(z+a_i)^{1-p}.
\]
At the upper point $z=\delta$ its apparent collision poles cancel.
To see this without a coefficientwise limit, choose a second, smaller
circle $|w|=\rho_0<\delta$ enclosing all $-a_i$. Then
\begin{equation}\label{hc:upper}
 U_m(\delta,\boldsymbol a)=\frac1{2\pi\ii}
  \int_{|w|=\rho_0}K_m(w,\boldsymbol a)
                         \log(\delta-w)\,\dd w.
\end{equation}
The logarithm is holomorphic inside this smaller circle. The displayed
integral is jointly holomorphic in the shifts. For $p\ge2$, the derivative
of the logarithm in the residue calculation is
$-(p-2)!/(\delta+a_i)^{p-1}$, giving exactly the sign in $U_m$.

At the lower point $z=-a_m+\eta$, the terms belonging to $i=m$ contain
only $\log\eta$ and negative powers of $\eta$; their cutoff constant
is zero. For $i>m$, ordinary substitution at $\eta=0$ is legitimate.
Subtracting these lower primitive values gives precisely the $m$th
summand in \eqref{hc:counter}. Consequently the remainder of the local
integral is
\[
 U_m(\delta,\boldsymbol a)
       +\int_{-a_m}^{\delta}W_m(z,\boldsymbol a)\,\dd z,
\]
which is jointly holomorphic, by \eqref{hc:Hermite} and
\eqref{hc:upper}. At zero shifts it is exactly the ambient finite part
of $K_m(z,\boldsymbol0)$ on $(0,\delta)$: all merged lower polar
primitive terms again have zero cutoff constant. Finally sum over $m$
and add the holomorphic all-regular and exterior integrals. The
zero-shift form of \eqref{hc:telescope} identifies their total with
$Q_{\boldsymbol r}$. This proves \eqref{hc:exact}.
\end{proof}

\begin{corollary}[Analytic collision arcs]\label{hc:arcs}
Let $a_i(\epsilon)$ be real analytic, vanish at zero, and be pairwise
distinct for sufficiently small $\epsilon>0$, with a fixed ordering
there. Then the complete expansion of
$\mathcal C_{\boldsymbol r}(\boldsymbol a(\epsilon))$ has the form
\[
 U(\epsilon)+V(\epsilon)\log\epsilon,
\]
where $U$ and $V$ are convergent Laurent germs with finite principal
parts. Its coefficient independent of $\epsilon$ and $\log\epsilon$ is
\begin{equation}\label{hc:arc-constant}
 \CT_{\epsilon}\mathcal C_{\boldsymbol r}(\boldsymbol a(\epsilon))
 =Q_{\boldsymbol r}
       +\CT_{\epsilon}\mathcal A_{\boldsymbol r}(\boldsymbol a(\epsilon)).
\end{equation}
Every term of the second constant is explicitly determined by finitely
many jets of the arcs and of the polygamma functions at $1$.
\end{corollary}

\begin{proof}
Each positive gap is
$a_i-a_m=b_{im}\epsilon^{\nu_{im}}(1+O(\epsilon))$, with
$b_{im}>0$ and positive integer $\nu_{im}$. Its logarithm is
$\nu_{im}\log\epsilon+\log b_{im}$ plus a convergent Taylor series.
Formula \eqref{hc:B} is meromorphic after composition with these arcs:
all possible poles come from a finite number of powers of nonzero
analytic gaps. Thus \eqref{hc:counter} has the stated form and its
constant needs only finitely many coefficients. The holomorphic
remainder contributes exactly its value at zero. No higher powers
of $\log\epsilon$ arise in this polygamma case.
\end{proof}

\subsection{A short formula at Stieltjes index zero}

Let $f=f_0=-\psi$ and $h(z)=f(1+z)$. If every $r_i=0$, all poles in
\eqref{hc:Km} are simple and \eqref{hc:counter} collapses to
\begin{equation}\label{hc:short}
 \boxed{\quad
 \mathcal A_{\boldsymbol0}(\boldsymbol a)=
 \sum_{m<i}\frac{\log(a_i-a_m)}{a_i-a_m}
   \prod_{k<m}h(a_k-a_i)
   \prod_{\substack{k>m\\k\ne i}}f(a_k-a_i).
 \quad}
\end{equation}
The negative arguments here are values of the meromorphic digamma
germ, not fractional parts. They are defined because the small gaps
are nonzero. Formula \eqref{hc:short} has only $d(d-1)/2$ summands.

For three factors at shifts $(0,a,b)$ it reads
\begin{equation}\label{hc:triple-exact}
 \mathcal A_3(0,a,b)=
 \frac{\log a}{a}f(b-a)
 +\frac{\log b}{b}f(a-b)
 +\frac{\log(b-a)}{b-a}h(-b).
\end{equation}
This exact identity is the key to the hierarchical limits below. For
fixed $a=\alpha\epsilon$, $b=\beta\epsilon$, expanding
\begin{equation}\label{hc:hseries}
 h(z)=\gamma+\sum_{k\ge1}(-1)^k\zeta(k+1)z^k
\end{equation}
recovers every displayed term through order $\epsilon^0$ in the
incoming fixed-rate triple theorem \cite{RepoMixed}. In particular,
for $(0,\epsilon,2\epsilon)$ the finite discrepancy is
$\zeta(2)\log2/2$, as already established there. The addition in
\eqref{hc:triple-exact} is the exact correction at arbitrary small
$a,b$, with a jointly holomorphic remainder, so rates may now coalesce.

Even a binary collision illustrates why the parameter matters. For
$d=2$, \eqref{hc:short} is $\log(a_2-a_1)/(a_2-a_1)$. Thus
\begin{equation}\label{hc:binary-curved}
 \CT_{\epsilon}\mathcal C_{0,0}(0,\epsilon+b\epsilon^2)
     =Q_{0,0}+b,
\end{equation}
because
$\log(\epsilon+b\epsilon^2)/(\epsilon+b\epsilon^2)
 =\epsilon^{-1}\log\epsilon-b\log\epsilon+b+O(\epsilon|\log\epsilon|)$.
This does not change the preceding binary theorem in its stated shift
coordinate. It records the effect of choosing a nonlinear collision arc.

\subsection{Two different hierarchies and their harmonic constants}

Set $Q_3=Q_{0,0,0}$, and write $H_n=\sum_{j=1}^n1/j$ for $n\ge1$.
In the formulas below a term containing $H_{u/v}$ is included only
when $v$ divides $u$; no value at a nonintegral index is needed.

\begin{theorem}[Oriented harmonic hierarchy laws]\label{hc:hierarchy}
Let $p>q\ge1$ be integers and $v=p-q$. The left nested collision obeys
\begin{equation}\label{hc:left}
 \boxed{\quad
 \CT_{\epsilon}\mathcal C_{0,0,0}(0,\epsilon^p,\epsilon^q)
   =Q_3-\mathbf1_{v\mid q}\,\gamma H_{q/v}.
 \quad}
\end{equation}
The right nested collision has the different value
\begin{equation}
 \boxed{\begin{aligned}
 \CT_{\epsilon}\mathcal C_{0,0,0}
           (0,\epsilon^q,\epsilon^q+\epsilon^p)
 =Q_3
 &+\mathbf1_{v\mid p+q}(-1)^{(p+q)/v}H_{(p+q)/v}
 \\[-2pt]
 &+\mathbf1_{v\mid q}\,\gamma(-1)^{q/v+1}H_{q/v}.
 \end{aligned}}\label{hc:right}
\end{equation}
In particular,
\begin{align}
 \CT_{\epsilon}\mathcal C_{0,0,0}(0,\epsilon^2,\epsilon)
     &=Q_3-\gamma,\label{hc:left-example}\\
 \CT_{\epsilon}\mathcal C_{0,0,0}(0,\epsilon,\epsilon+\epsilon^2)
     &=Q_3+\gamma-\frac{11}{6},\label{hc:right-example}\\
 \CT_{\epsilon}\mathcal C_{0,0,0}(0,\epsilon,\epsilon+\epsilon^3)
     &=Q_3+\frac32.\label{hc:right-second}
\end{align}
\end{theorem}

\begin{proof}
The two elementary, convergent generating functions needed are
\begin{equation}\label{hc:harmonic-generators}
 \frac{\log(1-z)}{1-z}=-\sum_{n\ge1}H_nz^n,
 \qquad
 \frac{\log(1+z)}{1+z}=\sum_{n\ge1}(-1)^{n+1}H_nz^n.
\end{equation}
The first follows by multiplying
$-\log(1-z)=\sum_{j\ge1}z^j/j$ by the geometric series; the second
is obtained by replacing $z$ with $-z$.

For the left nested configuration, substitute
$a=\epsilon^p$, $b=\epsilon^q$ into \eqref{hc:triple-exact}.
The first two logarithms are respectively $p\log\epsilon$ and
$q\log\epsilon$. They therefore make no contribution to the
log-free constant, regardless of the poles of their other factors.
The only possible contribution is
\[
 \epsilon^{-q}\frac{\log(1-\epsilon^v)}{1-\epsilon^v}
                       h(-\epsilon^q)
 =-\sum_{n\ge1}H_n\epsilon^{nv-q}
       \left\{\gamma+\sum_{k\ge1}\zeta(k+1)\epsilon^{qk}\right\}.
\]
For $k\ge1$, the exponent $nv-q+qk$ is strictly positive. The constant
comes solely from $k=0$ and $nv=q$, proving \eqref{hc:left} after
\eqref{hc:arc-constant} is used.

For the right nested configuration, put
$a=\epsilon^q$, $b=\epsilon^q+\epsilon^p$. Now the first and third
logarithms in \eqref{hc:triple-exact} are pure multiples of
$\log\epsilon$. The log-free part of the second term is
\begin{align*}
 &\frac{\log(1+\epsilon^v)}{1+\epsilon^v}
          \epsilon^{-q}f(-\epsilon^p)\\
 &\qquad=\left\{\sum_{n\ge1}(-1)^{n+1}H_n\epsilon^{nv}\right\}
       \left\{-\epsilon^{-p-q}+\gamma\epsilon^{-q}
                +\sum_{k\ge1}\zeta(k+1)\epsilon^{pk-q}\right\}.
\end{align*}
All terms in the last sum have positive exponents. The pole term
contributes only if $nv=p+q$, with coefficient $(-1)^nH_n$; the
constant regular term contributes only if $nv=q$, with coefficient
$\gamma(-1)^{n+1}H_n$. This proves \eqref{hc:right}. The examples use
$H_1=1$, $H_2=3/2$, and $H_3=11/6$.
\end{proof}

These two arrangements are not interchangeable. The original periodic
Stieltjes germ has a one-sided singularity, and reflection is not a
symmetry of its regular part. For example, subtracting
\eqref{hc:left-example} from \eqref{hc:right-example} eliminates the
unevaluated merged moment and gives the exact difference
$2\gamma-11/6$. The fixed-rate error estimate in the preceding work
alone would not establish either formula.

\subsection{A quartic example and an evaluated polygamma constant}

For linear rates $a_i=c_i\epsilon$, $c_1<\cdots<c_d$, define the
finite Laurent coefficient
\begin{equation}\label{hc:alpha}
 \alpha_{mi}(\boldsymbol c)=\frac1{c_i-c_m}
 [\epsilon^1]\left\{
   \prod_{k<m}h((c_k-c_i)\epsilon)
   \prod_{\substack{k>m\\k\ne i}}f((c_k-c_i)\epsilon)
 \right\}.
\end{equation}
Then the constant-order part of the complete counterterm is
\begin{equation}\label{hc:linear-constant}
 [\epsilon^0]\mathcal A_{\boldsymbol0}(\epsilon\boldsymbol c)
 =\sum_{m<i}\alpha_{mi}(\boldsymbol c)
               \{\log\epsilon+\log(c_i-c_m)\}.
\end{equation}
In particular, its log-free constant is the finite sum with
$\log(c_i-c_m)$. It requires only \eqref{hc:hseries} through degree
$d-2$. For rational rates every coefficient is a rational polynomial
in $\gamma,\zeta(2),\ldots,\zeta(d-1)$; its formal weight is $d-1$
when $\gamma$ has weight one and $\zeta(k)$ has weight $k$.
Indeed, a coefficient multiplying $z^j$ in either germ has weight
$j+1$, including the pole coefficient at $j=-1$; there are $d-2$
germ factors and their exponents sum to one in \eqref{hc:alpha}.
This is a statement about the produced expressions, not an assertion
of arithmetic independence.

For the quartic rates $(0,1,2,3)$ the six coefficients are
\begin{center}
\begin{tabular}{ccl}
\hline
$(m,i)$ & $c_i-c_m$ & $\alpha_{mi}$\\
\hline
$(1,2)$ & $1$ & $-3\gamma\zeta(2)+\tfrac92\zeta(3)$\\
$(1,3)$ & $2$ & $0$\\
$(1,4)$ & $3$ & $\gamma\zeta(2)-\tfrac32\zeta(3)$\\
$(2,3)$ & $1$ & $\gamma\zeta(2)+4\zeta(3)$\\
$(2,4)$ & $2$ & $2\gamma\zeta(2)-\tfrac92\zeta(3)$\\
$(3,4)$ & $1$ & $5\gamma\zeta(2)$\\
\hline
\end{tabular}
\end{center}
Consequently, with $Q_4=Q_{0,0,0,0}$,
\begin{align}
 \CT_{\epsilon}\mathcal C_{0,0,0,0}(0,\epsilon,2\epsilon,3\epsilon)
 =Q_4
 &+\gamma\zeta(2)(2\log2+\log3)\notag\\
 &-\frac32\zeta(3)(3\log2+\log3).
 \label{hc:quartic}
\end{align}
The coefficient of $\epsilon^0\log\epsilon$ is
$6\gamma\zeta(2)+5\zeta(3)/2$. A rate difference removes $Q_4$:
\begin{align}
 &\CT_{\epsilon}\mathcal C_{0,0,0,0}(0,\epsilon,2\epsilon,4\epsilon)
 -\CT_{\epsilon}\mathcal C_{0,0,0,0}(0,\epsilon,2\epsilon,3\epsilon)
 \notag\\
 &\qquad=\gamma\zeta(2)\left(\frac72\log2+\log3\right)
       -\frac{\zeta(3)}6(\log2+7\log3).
 \label{hc:quartic-difference}
\end{align}

Repeated poles add the rational primitive terms in
\eqref{hc:counter}; retaining only its logarithmic terms would give
an incorrect finite constant. For example,
\begin{equation}\label{hc:middle-derivative}
 \boxed{\quad
 \CT_{\epsilon}\FP_x\int_0^1
   \gamma_0(x)\gamma_0'(\{x+\epsilon\})
                         \gamma_0(\{x+2\epsilon\})\,\dd x
 =2\gamma\zeta(2)-\zeta(3)\log2.
 \quad}
\end{equation}
To check every term, for $\boldsymbol r=(0,1,0)$ and rates $(0,1,2)$,
the three pairs in \eqref{hc:counter} have constant-order contributions
\[
 (m,i)=(1,2):\ -3\zeta(3)\log\epsilon+\zeta(3),\qquad
 (1,3):\ -\zeta(3)(\log\epsilon+\log2),\qquad
 (2,3):\ 4\zeta(3)\log\epsilon.
\]
The nonlogarithmic $\zeta(3)$ in the first pair comes from its double
pole. Hence the finite discrepancy from the merged moment is
$\zeta(3)(1-\log2)$. The merged value follows directly from the
fundamental theorem of calculus before taking the cutoff constant:
\[
 Q_{0,1,0}=\frac13\left\{\gamma^3-
                 \CT_{x\downarrow0}\gamma_0(x)^3\right\}
          =2\gamma\zeta(2)-\zeta(3),
\]
since the constant coefficient of
$(x^{-1}+\gamma-\zeta(2)x+\zeta(3)x^2+\cdots)^3$ is
$\gamma^3-6\gamma\zeta(2)+3\zeta(3)$. Adding the discrepancy proves
\eqref{hc:middle-derivative} without an unevaluated global moment.

\subsection{Verification and scope of the result}

The supplied \src{code/verify\_collisions.py} checks both harmonic
hierarchy laws by exact coefficient extraction for every
$1\le q<p\le12$ and verifies the quartic expressions. Its numerical
calculation independently integrates the original periodic digamma
products, using local Taylor subtraction at each separated seam. The
merged moments are obtained from a local Laurent expansion and ordinary
quadrature away from zero. The collision formula itself is not used
in either quadrature. At 70-digit working precision, selected residuals
$\mathcal C-\mathcal A-Q$ are
\begin{center}
\begin{tabular}{lrr}
\hline
Configuration & $\epsilon=10^{-2}$ & $\epsilon=10^{-3}$\\
\hline
$(0,\epsilon,2\epsilon,3\epsilon)$
 & $-1.468167310\cdot10^{-1}$ & $-1.501685165\cdot10^{-2}$\\
$(0,\epsilon^2,\epsilon)$
 & $-1.288859769\cdot10^{-2}$ & $-1.210723679\cdot10^{-3}$\\
$(0,\epsilon^3,\epsilon^2)$
 & $-1.202913331\cdot10^{-4}$ & $-1.202065558\cdot10^{-6}$\\
$(0,\epsilon,\epsilon+\epsilon^2)$
 & $-1.268176013\cdot10^{-2}$ & $-1.208676388\cdot10^{-3}$\\
\hline
\end{tabular}
\end{center}
These residuals are consistent with a holomorphic remainder, whose
first change is of the order of the largest shift. They are not claimed
to be interval certificates. A separate repeated-pole quadrature in
\src{code/verify\_polygamma\_collisions.py} checks
\eqref{hc:middle-derivative} using second-order local subtraction;
the corresponding residuals are approximately $0.02934754$ and
$0.00224172$. Exact coefficient generation for that example and
additional derivative placements is supplied in
\src{code/derive\_polygamma\_collisions.py}.

The proof of unrestricted multiplicity, argument-derivative order, and
analytic collision rates is the telescoping and contour argument above.
The numerical evidence checks its implementation on specified cases.
The prior fixed-rate triple discrepancy and fixed-collision spectral
completion remain credited to \cite{RepoMixed}. Positive Stieltjes
indices introduce logarithmic singular germs rather than meromorphic
ones; a corresponding exact geometric formula would require an
additional argument and is not asserted here. The unresolved extension
is precisely formulated in the research agenda.

% END sections/05_collisions.tex

% BEGIN sections/06_audit_verification.tex
\section{Source audit, proposed qualifications, and verification}
\label{sec:audit}

\subsection{A fixed source record}

The main source baseline is the ProveIt commit whose full identifier is
\begin{center}
\small\src{921b58f1099568bfefa28a556470e4b29e94eb3b}.
\end{center}
It contains the manuscript directory supplied for this task and four
incoming archives dated 10--11 October 2026: the bilinear
Mellin--Hurwitz, cubic harmonic mixed-order, exact resonant, and
independent-order reports \cite{RepoMain,RepoBilinear,RepoCubic,
RepoExactResonant,RepoIndependent}. The earlier mixed spectral report,
which states several questions answered here, was inspected at the
historical revision
\begin{center}
\small\src{5a790187c8e186e41e2b990b4941cb7a1a3c7b6b}.
\end{center}
Its archived source is cited separately as \cite{RepoMixed}; earlier
exact identities used in the shift calculation are \cite{RepoExact}.
This distinction matters because an incoming archive can be integrated
and removed from a later incoming directory without invalidating its
historical role.

The audit covered the definitions, proofs, status statements, and
questions directly used in this continuation, together with the current
Gaussian conjecture labels. It did not certify every line of the entire
multi-volume manuscript. No new false theorem was identified in the
selected source material. The qualifications below state how the new
results must be recorded; they are not allegations that the source
already contains each invalid inference.

\subsection{Status changes justified by the proofs}

\paragraph{The shift-boundary question is answered.}
The mixed spectral report asks for the actual germ at $a=-r$ and the
exception at nonpositive spectral orders. Theorems~\ref{hsh:germ},
\ref{hsh:boundary}, \ref{hsh:radius}, and \ref{hsh:rank} answer it.
The existing half-plane theorem remains correct, but a status summary
can now give the exact first singularity and exact Taylor radius.
Transport itself should be attributed to the broader relative ordered
Hurwitz character: the exact identification is
\eqref{hsh:relativecharacter}. Treating the present family as unrelated
to that source would duplicate an established construction.

\paragraph{The weight-six and weight-eight separation question is answered
within its specified family.}
Equations~\eqref{bell:sixrref} and \eqref{bell:eightrref} give the
requested rational elimination. Theorem~\ref{bell:rank} supplies the
general rank, while Theorem~\ref{bell:image} determines the full
summation kernel. The individual nonlinear coordinate moments at
weights six and eight lie outside this formal mixed-Bell span.
This does not rule out evaluating them using another parameter family,
an argument derivative, or an independent functional identity.

\paragraph{The compact harmonic-power target has a constructive solution.}
Theorem~\ref{hpnew:generator} and
Corollary~\ref{hpnew:threerows} answer the complete minus-two sequence
requested in Question~10 of \cite{RepoExactResonant}; they also give
every further centered even moment. Theorem~\ref{hpnew:uniformspan}
adds the stronger conclusion that one finite Euler-sum span works for
all moment orders at a fixed power. A closed bivariate expression and
minimal arithmetic coordinates remain open.
Corollary~\ref{hpnew:oddorders} supplies a sufficient generalized-harmonic
parity class toward Question~11 and a finite formula for every single
positive odd harmonic order; it does not finish the requested necessity
classification.

\paragraph{Nested collision constants are now explicit in the polygamma
sector.}
The exact formula of Theorem~\ref{hc:main} supplies a jointly
holomorphic remainder for any number of factors and any nonnegative
argument-derivative orders. Theorem~\ref{hc:hierarchy} answers the
named nested triple paths and all their monomial-rate analogues.
The extension to positive Stieltjes indices and a theorem about
composition of renormalized merging operations are still separate
questions. They should not be marked solved by the present result.

\subsection{Qualifications that prevent incorrect extensions}

\paragraph{Principal continuation is a statement about a branch.}
The principal family is holomorphic on the slit plane in
\eqref{hsh:slit}. Earlier removable shifts need not remain removable
on every sheet. A loop around $a=-1$ changes $D_1$ by
\[
 \bigl(e^{-2\pi\ii s}-1\bigr)\frac{(a+1)^{-s}}a.
\]
The added term generally has a pole at $a=0$. A claim of unrestricted
removability on the universal continuation would therefore need
correction. The source's valid half-plane result makes no such claim.

\paragraph{Polynomial values do not imply polynomial spectral derivatives.}
The finite Stirling formula at $s=-m$ has rank $m+1$ in depth.
Differentiating the analytic family introduces
$z^m(-\log z)^k$ at the first boundary. For every $k\ge1$ it has
finite shift Taylor radius and supplies independent depth functions.
It cannot be obtained by differentiating a formula whose label $m$
has only been specified at integers. Similarly, the first spectral
derivative of the raw harmonic-power family retains a global Mellin
remainder that is absent from its regular value.

\paragraph{A removable restriction need not be a jointly removable germ.}
The auxiliary $G_m$ in Section~\ref{hpnew:section} obeys $G_m(0)=0$,
whereas separately specializing the uncentered family at exponent zero
gives $\zeta(-m)$. The explicit term in
Remark~\ref{hpnew:orderwarning} proves the discrepancy. Its polar
hyperplane is not removable merely because its numerator vanishes at
one point. In contrast, the centered family has no polar hyperplane
through $(-2r,0)$, so the claimed coefficient extraction is valid.

\paragraph{A collision constant depends on the full primitive.}
The terms of order $p\ge2$ in \eqref{hc:counter} are rational
primitive terms, in addition to logarithms. They are essential:
the derivative example \eqref{hc:middle-derivative} would be wrong
without them. Even a curved binary gap $\epsilon+b\epsilon^2$
changes the raw constant by $b$, as
\eqref{hc:binary-curved} shows. Normalizing only a leading pole or
only a fixed-rate approximation does not determine the required
constant along a nested path.

\paragraph{Formal independence and arithmetic independence are distinct.}
The Bell rank treats the $L_k$ as independent indeterminates, and the
shift rank treats analytic functions of $a$ over the constants.
The Euler sums in $\mathcal E_{p-1}$ form a proved spanning list.
None of these assertions gives a lower bound for the arithmetic
dimension of their numerical values. The exact matrices are suitable
certificates for the specified algebraic problem and no larger one.

\paragraph{The retained Gaussian conjectures remain conjectures.}
The pinned manuscript labels the current $S_6$ candidate
\src{cycloquot:conj:S6} in
\src{chapters/04-cyclotomic-quotients.tex}, and the revised $S_8$
candidate \src{s8new:conj:S8} in
\src{chapters/04-S8-candidate.tex} \cite{RepoMain}.
This article proves neither numerical period identity. The rejected
older $S_8$ vector is already distinguished in the current source.
No new integer-relation fit or conjectural coefficient vector is
offered here. The remaining task is a new exact functional relation
with its analytic endpoint evaluation.

\subsection{Reproducible exact checks and independent diagnostics}

Every primary analytic result has a proof in the preceding sections.
The following additional checks are included as readable Python source
and JSON results. They require only Python, SymPy, and mpmath.
The driver \src{code/run_all.py} runs the checks in a fixed order.

\begin{longtable}{@{}p{0.24\textwidth}p{0.69\textwidth}@{}}
\toprule
Family & Recorded verification\\
\midrule
\endhead
Shift germs & 85 exact core checks and 50 numerical comparisons:
Newton versus defining sums, finite transport, Cauchy-extracted
boundary constants, spectral jets, and a doubled cutoff.
The primitive supplement adds 12 exact derivative and partial-fraction
checks.\\[5pt]
Mixed Bell moments & Independent coefficient constructions and ranks
for weights $2$ through $18$; exact elimination and annihilators at
weights $6,8$; the quartic zero row; and a deliberately corrupted
coefficient that is correctly rejected. Nineteen mixed rows at five
real or complex parameter choices are checked numerically, plus the
independent harmonic tail-square identity.\\[5pt]
Harmonic powers & Seven exact rational generator residuals;
26 Bernoulli identities through moment order $r=12$;
15 centered values at powers $1$--$5$ and moment orders $0$--$2$;
and six full exponent-generator comparisons at two exponent values.
The centered computations use two truncation lengths.\\[5pt]
Generalized harmonics & Nine exact centered Faulhaber identities and
six independent values for harmonic orders $q=3,5$ at moment orders
$r=0,1,2$, also using two subtraction truncations.\\[5pt]
Collisions & 132 exact oriented hierarchy assertions for
$1\le q<p\le12$, exact quartic coefficients and rate differences,
and the derivative constant. Independent local-subtraction quadrature
is used for 12 digamma configurations and two polygamma configurations.
The corrected remainder tends to the merged value as predicted.\\
\bottomrule
\end{longtable}

The largest scaled residual in the core shift comparisons is about
$1.76\cdot10^{-23}$ and is controlled by the chosen Newton cutoff.
The Bell mixed-row comparisons are within
$2.70\cdot10^{-65}$ in the recorded scaled metric; the tail-square
comparison has absolute residual below $8.30\cdot10^{-75}$.
For the 15 centered harmonic-power values, the largest absolute
discrepancy is below $6.49\cdot10^{-45}$; the independent truncations
differ by at most $1.29\cdot10^{-32}$. These are recorded diagnostic
values, not rigorous error estimates. The precision, truncation
parameters, and individual residuals are retained in the result files.

The collision diagnostics test the exact correction
$\mathcal C-\mathcal A$ against its proved holomorphic limit.
They do not assert that this difference equals its limit at a nonzero
shift: the remaining Taylor terms are expected. This makes their
numerical interpretation different from the residual of an exact
finite-parameter Bell identity.

\subsection{Integration order}

The source sections can be integrated independently, retaining their
label prefixes \src{hsh:}, \src{bell:}, \src{hpnew:}, and \src{hc:}.
First reconcile the relative-character identification, then add the
shift singularities and primitives; the Bell and harmonic-power
sections can follow their existing antecedents; finally insert the
collision theorem before the nested-rate examples. The plain-text
integration note records exact source dependencies and suggested
status replacements. The complete standalone TeX and its compiled PDF
are included, alongside the modular source and verification programs.

% END sections/06_audit_verification.tex

% BEGIN sections/07_research.tex
\section{Further research questions and exact targets}
\label{sec:research}

The questions below use the new results as starting points. Each asks
for a functional identity, a finite reduction, or a specified algebraic
classification. None is asserted as a proved conclusion of this article.

\subsection{Harmonic powers and generalized harmonic numerators}

\begin{question}[A single generating expression in both orders]
The triangular generator determines every $V_r(u)$ using $2r+1$
steps. Find a useful closed expression for
\[
 \sum_{r\ge0}V_r(u)\frac{t^{2r}}{(2r)!}
\]
in a stated domain, with its singular terms in $t$ separated and its
continuation in $u$ controlled. A satisfactory expression should
recover both the negative odd exponent poles and the uniform finite
Euler-sum span without first expanding all the triangular rows.
The centered Mellin kernel and its even local coefficients give a
concrete route. An identity of formal coefficient series should be
distinguished from a normally convergent two-variable identity.
\end{question}

\begin{question}[Minimal coordinates at a fixed harmonic power]
For each $p$, determine the span of all $F_p(-2r,1/2)$ inside the
finite displayed space $\mathcal E_{p-1}$. Compute rational
coefficient matrices for successive $r$ and reduce them modulo an
explicitly stated set of proved multiple-zeta relations. The first
targets are $p=6,7,8$, where the generating theorem is finite but
individual Euler sums begin to interact nontrivially. The desired
output is a minimal presentation relative to those proved relations,
not an unsupported assertion of arithmetic independence.
\end{question}

\begin{question}[A criterion for regular generalized-harmonic products]
Classify polynomials in $H_n^{(q)}$ whose centered Dirichlet series
are holomorphic at every nonpositive even spectral point. The
asymptotic coefficients in inverse powers of $n+1/2$ give necessary
local tests. Determine which combinations involving even harmonic
orders cancel all obstructing coefficients, and express the answer
by a generating identity rather than an expanding list of examples.
This continues Question~11 of \cite{RepoExactResonant} and should
include both sufficient structural families and a clearly specified
necessity statement.
\end{question}

\begin{question}[Exact cancellation of the first global spectral term]
Find finite combinations of the centered harmonic-power families
whose first spectral derivatives at $s=-2r$ reduce to a finite
Euler-sum and Stieltjes expression. The regular values alone cannot
settle this question: the holomorphic Mellin remainder contributes
after differentiation. A proof should first exhibit an exact
identity among the Mellin kernels that cancels that remainder, then
differentiate the completed family. A first target is a relation
among powers $p\le4$ at $s=0$ or $s=-2$.
\end{question}

\subsection{Shift singularities, monodromy, and primitives}

\begin{question}[The full relative-character singularity pattern]
Extend the explicit scalar amplitudes in Theorem~\ref{hsh:germ}
from words $(s,\{1\}^r)$ to arbitrary independent complex orders
in the relative Hurwitz character of \cite{RepoIndependent}.
At each negative shift, determine a finite set of local power--log
germs and their exact coefficients. Then identify which words have
removable earlier shifts and which exceptional order loci reduce the
functional rank. The source transport is available; the challenge
is organizing the cancellations without hiding them in a recursively
defined remainder.
\end{question}

\begin{question}[A monodromy representation with explicit generators]
The local formula gives the monodromy around each negative integer.
Describe the action of finite products of these loops on the span
of relative-character specializations, including the rational poles
created on nonprincipal sheets. Determine whether a finite-depth
filtration is preserved and compute its associated graded action.
This would give a global continuation statement that retains the
branch information responsible for the exact-radius theorem.
\end{question}

\begin{question}[Global normalized primitives and shift transport]
The local primitive in Section~\ref{hsh:primitives} is explicit in
Lerch functions, with finite polylogarithmic formulas for spectral
jets at zero. Construct global primitives normalized at a positive
shift, and determine how their constants transform under
$a\mapsto a+1$ and under monodromy. A useful next target is a
closed difference equation for the first primitive of $D_1$ and its
first two spectral derivatives, retaining every logarithmic contact
term. Repeated primitives may connect to higher Gamma functions,
but any such identification must fix the same base-point constants.
\end{question}

\subsection{Polygamma moments beyond the two-parameter rows}

\begin{question}[Additional generators that isolate the missing moments]
The exact Bell matrices identify the missing formal coordinates at
weights six and eight. Find an additional Gamma or beta parameter
identity whose mixed derivatives contribute a row outside that span.
The immediate target at weight six is an individual value of
\[
 \sum_{n\ge0}(n+a/2)L_2(n;a)^3
\]
for generic $a$, since the two combinations in
\eqref{bell:sixidentity1}--\eqref{bell:sixidentity2} then separate
the other cubic-weight products. At weight eight, choose candidate
generators by pairing their coefficient rows with the explicit
annihilators in \eqref{bell:eightann} before carrying out numerical
experiments. This gives an exact test of whether a proposed family
adds information.
\end{question}

\begin{question}[A hierarchy of generalized-harmonic tail identities]
Specialize the full summation kernel in
\eqref{bell:kernelbasis} at integral and half-integral $a$.
Classify the resulting multiple-zeta or colored multiple-polylogarithm
identities at successive weights and compare their span with stated
shuffle and stuffle relations. The weight-seven tail-square
evaluation is the first compact example. The target is a finite
family valid at arbitrary weight, together with a transparent map
from Bell rows to convergent harmonic sums, rather than a list of
numerically fitted reductions.
\end{question}

\subsection{Collision constants and higher Stieltjes indices}

\begin{question}[Positive Stieltjes indices with moving gaps]
Extend Theorem~\ref{hc:main} to products of
$\gamma_{m_i}^{(r_i)}(\{x+a_i\})$ with some $m_i>0$.
Their singular pieces contain logarithms, so a meromorphic
partial-fraction decomposition alone is insufficient. Introduce
independent local exponents, integrate before differentiating, and
seek a jointly holomorphic remainder after subtraction of the full
parameter-dependent primitives. The first target is three factors
with indices $(1,0,0)$ along both nested paths of
Theorem~\ref{hc:hierarchy}. The statement should specify the
allowed logarithm branches and every intermediate normalization.
\end{question}

\begin{question}[Composition of renormalized collisions]
The exact correction isolates a holomorphic remainder at the full
collision. Can one define partial-cluster corrections so that merging
two clusters and then a third produces the same normalized result as
the other parenthesization? Raw constant extraction is already
nonassociative in the relevant sense, as the oriented harmonic laws
show. Determine the exact defect and whether a local correction
removes it. Start with four digamma factors and two nested clusters;
the primitive coefficients in \eqref{hc:B} give finite algebraic
data for testing the compatibility equations.
\end{question}

\begin{question}[Global multiple kernels after local collision subtraction]
The present theorem isolates all local gap dependence, leaving the
merged constant $Q_{\boldsymbol r}$. Identify permutation or
argument-derivative combinations for which this global constant
reduces to endpoint data, extending the evaluated example
\eqref{hc:middle-derivative}. For the remaining combinations,
give an explicit finite decomposition into convergent multiple
Tornheim-type kernels. Begin with four factors and small total
argument-derivative order; distinguish a local completion theorem
from an evaluation of the surviving global period.
\end{question}

\subsection{The unresolved Gaussian candidates}

\begin{question}[An exact functional bridge for the frozen $S_6$ and $S_8$]
Retain the current source coefficient vectors and the known formal
nonmembership results \cite{RepoMain,RepoExactResonant}. Seek an
additional iterated-integral, distribution, or functional relation
outside the already excluded finite relation families. Any claimed
proof should supply both an exact rational certificate and the
analytic identity represented by each row, including endpoint
regularization. A new relation search inside an unchanged excluded
span cannot settle the conjecture. The Bell and harmonic identities
here provide exact auxiliary periods, but no bridge from them to
the Gaussian targets has been established.
\end{question}

The most immediate extensions are the fixed-power coordinate matrices,
the first positive-Stieltjes-index nested collision, and a new Gamma
generator selected using the weight-six annihilator. Each begins from
a finite, explicit calculation and asks for an exact functional
identity. The broader monodromy and Gaussian questions require
additional structural input.

% END sections/07_research.tex

% BEGIN references.tex
\begin{thebibliography}{99}
\addcontentsline{toc}{section}{References}

\bibitem{RepoMain}
Vladimir Reshetnikov, \emph{ProveIt}, polylogarithms manuscript and
incoming research directory, repository revision
\src{921b58f1099568bfefa28a556470e4b29e94eb3b}.
\href{https://github.com/VladimirReshetnikov/ProveIt/tree/921b58f1099568bfefa28a556470e4b29e94eb3b/Analysis/Polylogarithms/docs/manuscript}{Pinned manuscript};
\href{https://github.com/VladimirReshetnikov/ProveIt/tree/921b58f1099568bfefa28a556470e4b29e94eb3b/docs/incoming}{pinned incoming directory}.
Inspected 11 October 2026.

\bibitem{RepoMixed}
ProveIt research archive,
\emph{Mixed Spectral Identities},
\nolinkurl{ProveIt_Mixed_Spectral_Identities_2026-10-10.zip},
historical revision \src{5a790187c8e1}.
Especially \src{sections/02_harmonic.tex},
\src{sections/03_dougall.tex},
\src{sections/05_collisions.tex}, and
\src{sections/06_audit_research.tex}.
\href{https://github.com/VladimirReshetnikov/ProveIt/blob/5a790187c8e186e41e2b990b4941cb7a1a3c7b6b/docs/incoming/ProveIt_Mixed_Spectral_Identities_2026-10-10.zip}{Pinned archive}.

\bibitem{RepoExact}
ProveIt research archive,
\emph{Exact Identities},
\nolinkurl{ProveIt_Exact_Identities_2026-10-10.zip},
historical revision \src{5a790187c8e1}.
\href{https://github.com/VladimirReshetnikov/ProveIt/blob/5a790187c8e186e41e2b990b4941cb7a1a3c7b6b/docs/incoming/ProveIt_Exact_Identities_2026-10-10.zip}{Pinned archive}.

\bibitem{RepoIndependent}
ProveIt research archive,
\emph{Independent Orders},
\nolinkurl{ProveIt_Independent_Orders_2026-10-11.zip},
revision \src{921b58f10995}.
Especially \src{sections/02_hurwitz.tex}, labels
\src{eht:rec}, \src{eht:lower-shift}, and \src{eht:ones-Gamma}.
\href{https://github.com/VladimirReshetnikov/ProveIt/blob/921b58f1099568bfefa28a556470e4b29e94eb3b/docs/incoming/ProveIt_Independent_Orders_2026-10-11.zip}{Pinned archive}.

\bibitem{RepoExactResonant}
ProveIt research archive,
\emph{Exact Resonant Identities},
\nolinkurl{ProveIt_Exact_Resonant_Identities_2026-10-11.zip},
revision \src{921b58f10995}.
Especially the raw harmonic-power section and
\src{sections/07_research.tex}, Questions 10--12.
\href{https://github.com/VladimirReshetnikov/ProveIt/blob/921b58f1099568bfefa28a556470e4b29e94eb3b/docs/incoming/ProveIt_Exact_Resonant_Identities_2026-10-11.zip}{Pinned archive}.

\bibitem{RepoCubic}
ProveIt research archive,
\emph{Cubic Harmonic Mixed Orders},
\nolinkurl{ProveIt_Cubic_Harmonic_Mixed_Orders_2026-10-11.zip},
revision \src{921b58f10995}.
\href{https://github.com/VladimirReshetnikov/ProveIt/blob/921b58f1099568bfefa28a556470e4b29e94eb3b/docs/incoming/ProveIt_Cubic_Harmonic_Mixed_Orders_2026-10-11.zip}{Pinned archive}.

\bibitem{RepoBilinear}
ProveIt research archive,
\emph{Bilinear Mellin--Hurwitz},
\nolinkurl{ProveIt_Bilinear_Mellin_Hurwitz_2026-10-10.zip},
revision \src{921b58f10995}.
\href{https://github.com/VladimirReshetnikov/ProveIt/blob/921b58f1099568bfefa28a556470e4b29e94eb3b/docs/incoming/ProveIt_Bilinear_Mellin_Hurwitz_2026-10-10.zip}{Pinned archive}.

\bibitem{ArakawaKaneko}
T. Arakawa and M. Kaneko,
Multiple zeta values, poly-Bernoulli numbers, and related zeta functions,
\emph{Nagoya Mathematical Journal} \textbf{153} (1999), 189--209.
\href{https://doi.org/10.1017/S0027763000006954}{doi:10.1017/S0027763000006954}.
\href{https://www2.math.kyushu-u.ac.jp/~mkaneko/papers/18MZV_pB_Nagoya.pdf}{Author-hosted text}.

\bibitem{KanekoTsumura}
M. Kaneko and H. Tsumura,
Multi-poly-Bernoulli numbers and related zeta functions,
\emph{Nagoya Mathematical Journal} \textbf{232} (2018), 19--54.
\href{https://doi.org/10.1017/nmj.2017.16}{doi:10.1017/nmj.2017.16}.
\href{https://www2.math.kyushu-u.ac.jp/~mkaneko/papers/56Kaneko_Tsumura.pdf}{Author-hosted text}.
The increasing-index convention there is reversed when compared with
the decreasing-index convention stated in this article.

\bibitem{DLMFHurwitz}
NIST Digital Library of Mathematical Functions,
\S25.11, Hurwitz zeta function.
\url{https://dlmf.nist.gov/25.11}.
Especially the shift recurrence, polygamma relation,
nonpositive integer values, and derivatives.

\bibitem{DLMFPolygamma}
NIST Digital Library of Mathematical Functions,
\S5.15, Polygamma functions.
\url{https://dlmf.nist.gov/5.15}.

\bibitem{DLMFGammaAsymptotic}
NIST Digital Library of Mathematical Functions,
\S5.11, Gamma-function asymptotic expansions.
\url{https://dlmf.nist.gov/5.11}.

\bibitem{DLMFDougall}
NIST Digital Library of Mathematical Functions,
\S16.4, Argument unity, equation 16.4.9,
Dougall's very-well-poised ${}_5F_4(1)$ summation.
\url{https://dlmf.nist.gov/16.4.E9}.

\bibitem{DLMFPolylog}
NIST Digital Library of Mathematical Functions,
\S25.12, Polylogarithms.
\url{https://dlmf.nist.gov/25.12}.

\bibitem{DLMFLerch}
NIST Digital Library of Mathematical Functions,
\S25.14, Lerch's transcendent.
\url{https://dlmf.nist.gov/25.14}.

\end{thebibliography}

% END references.tex

\end{document}
